\documentclass[11pt]{article}
\usepackage[a4paper,margin=19mm]{geometry}
\usepackage[no-math]{fontspec}
\setmainfont{Latin Modern Roman}[SmallCapsFont=Latin Modern Roman Caps]
\usepackage{amsmath,amssymb,amsthm,mathtools,amscd,verbatim,xcolor,graphicx,hyperref}
\hypersetup{hidelinks}
\numberwithin{equation}{section}
\newtheorem{theorem}{Theorem}[section]
\newtheorem{lemma}[theorem]{Lemma}
\newtheorem{proposition}[theorem]{Proposition}
\newtheorem{corollary}[theorem]{Corollary}
\newtheorem{claim}[theorem]{Claim}
\newtheorem{fact}[theorem]{Fact}
\newtheorem{conjecture}[theorem]{Conjecture}
\theoremstyle{definition}
\newtheorem{definition}[theorem]{Definition}
\newtheorem{example}[theorem]{Example}
\newtheorem{remark}[theorem]{Remark}
\newtheorem{question}[theorem]{Question}
\theoremstyle{plain}
\newcommand{\C}{\mathbb C}
\newcommand{\R}{\mathbb R}
\newcommand{\Q}{\mathbb Q}
\newcommand{\Z}{\mathbb Z}
\newcommand{\N}{\mathbb N}
\newcommand{\eee}{\mathrm e}
\newcommand{\cclass}[1]{\textsf{#1}}
\newcommand{\cproblem}[1]{\textsc{#1}}
\newcommand{\ie}{i.\,e.}
\newcommand{\Ie}{I.\,e.}
\newcommand{\eg}{e.\,g.}
\newcommand{\Eg}{E.\,g.}
\newcommand{\phd}{Ph.\,D.}
\newcommand{\msc}{M.\,S.}
\newcommand{\bsc}{B.\,S.}
\newcommand{\expref}[2]{\hyperref[#2]{#1~\ref{#2}}}
\newcommand{\expeqref}[2]{\hyperref[#2]{#1~\eqref{#2}}}
\newcommand{\exprefsilent}[2]{\hyperref[#2]{#1}}
\newcommand{\epfmtdoi}[1]{\href{https://doi.org/#1}{doi:#1}}
\emergencystretch=3em
\usepackage{etoolbox,bm}
\newcommand{\epfmt}[2]{\begingroup\def\TocNativeBibResult{\PackageError{toc-native-compat}{Unsupported bibliography type/record #1/#2}{Only exact source-observed finite records are allowed.}}\TocNativeBibResult\endgroup}

%% Theory of Computing
%% File: v021a008.tex
%% Publication date: October 28, 2025
%% Authors: Gil Cohen, Dean Doron, and Shahar Samocha
%
%
%
%
%

%
%

%
%

%
%

%
%

%
%

%
%

%
%
%
%

%

%
%
%
%
%
%
%
%
%
%
%
%
%
%
%
%
%
%
%
%
%
%

%



%
%
%
   %

    %
    %
    %
    %
    %
    %
    %
    %
    %
    %
    %
    %
    %
    %
    %
    %
    %
    %
    %
    %
    %
    %
    %
    %
    %
    %
    %
    %
    %

%
%
   %

%
%
%
%
%
%

%
%
%

\makeatletter
\newcommand{\ccc}{\mathcal{C}}
\newcommand{\ext}{\mathsf{Ext}}
\newcommand{\adv}{A}
\newcommand{\advfam}{\mathcal{A}}
\newcommand{\advpair}{\mathcal{T}}
\newcommand{\eps}{\varepsilon}
\newcommand{\perfam}{\Pi}
\newcommand{\nc}{\mathcal{X}}
\newcommand{\me}{{H_{\infty}}}
\newcommand{\set}[1]{\left\{{#1}\right\}}
\newcommand{\B}{\set{0,1}}
\newcommand{\E}{\mathop{\bf E\/}}
\newcommand{\abs}[1]{\left\lvert #1\right\rvert}
\newcommand{\advfampart}{\mathcal{F}}
\newcommand{\nofp}{\mathcal{N}}
\newcommand{\minent}[1]{H_{\infty}(#1)}
\newcommand{\ane}{\mathrm{NE}}
\newcommand{\lseed}{Y_{[2:d]}}
\def\poly{{\mathrm{poly}}}
\def\polylog{{\mathrm{polylog}}}
\def\SD{\ensuremath{\mathsf{SD}}}	
\def\MM{{\mathcal{M}}}
\usepackage{bm}
\makeatother


\begin{document}
\begin{center}{\Large Stable item 748: Existence of multi-seed-protecting extractors against noncolluding entropy-preserving adversaries}\end{center}
\paragraph{Source and attribution.}
Gil Cohen, Dean Doron, and Shahar Samocha, \emph{Seed-Protecting Extractors}. Theory of Computing 21 (2025), Article 8, publisher source publication October 28, 2025; DOI 10.4086/toc.2025.v021a008. \href{https://theoryofcomputing.org/articles/v021a008/}{Publisher article}. Original human statement/proof/context: \href{https://creativecommons.org/licenses/by/3.0/}{CC BY3.0}.
\paragraph{Selected problem and source scope.}
Prove only the authors' Theorem 6.1 (original label \detokenize{thm:non-explicit}) below. The full original mathematical article and bibliography are preserved as conservative closed context. The exact selected theorem is an asymptotic existence result with a displayed log(n-k) seed bound, so its logarithm has the source expression's n greater than k domain. The final proof uses n-k+2 in the sufficient seed-size inequality; no new exact n=k endpoint is asserted. The extractor and noncolluding-adversary definitions remain those of the original article, with both distinguishability terms and the full collision-graph/martingale core included. The sibling seed-protection-to-nonmalleability branch remains unqualified context due to its documented pool-cardinality question.  Other equivalent formulations, consequences and supporting lemmas are uncounted.
\paragraph{Difficulty (editorial): H3.}
The existence proof has to control random extractor evaluations that overlap across several noncolluding tampering functions, so its main concentration argument cannot assume independent samples. The authors build a bounded-degree collision graph, partition its seeds into independent classes, and establish the conditional mean-zero property needed for separate martingale bounds on both the joint and product-distribution distinguishers. Their complete human core then performs the simultaneous source, adversary and test-function union bound with explicit entropy and seed-length costs; the difficulty is this new dependency-control construction rather than a routine random-function count.
\paragraph{Editorial source boundary.}
Author mathematical TeX and ordering are unchanged. Only journal front matter is replaced by this wrapper; exact authored macros are flattened and separately preserved, and finite source-observed display aliases are defined in the standalone preamble. There are no original figures or external mathematical assets in this paper.  No generated proof, mathematical repair or independent mathematical refereeing is supplied. 
\clearpage

%
%
%


\section{Introduction}


Informally, a \emph{seeded extractor} is a function that ``purifies'' defective randomness using few fresh  random bits. A defective random source is modelled by a distribution $X$ that has some
%
 lower bound on its \emph{min-entropy}. A random variable $X$ is said to have min-entropy $k$, denoted $\me(X) \ge k$, if for every $x$, $\Pr[X = x] \le 2^{-k}$.  When $X$ is supported over $n$-bit strings, we call $X$ an $(n,k)$-source. A function $\ext \colon \{0,1\}^n \times \{0,1\}^d \to \{0,1\}^m$ is a \emph{$(k,\eps)$-seeded extractor}~\cite{NZ96} if for every $(n,k)$-source $X$ it holds that $\ext(X,Y)$ is $\eps$-close, in statistical distance, to the uniform distribution over $m$-bit string. Here, $Y$ is a random variable, independent of $X$, that is uniformly distributed over $d$-bit strings. We write this as $\ext(X,Y) \approx_\eps U_m$. Informally, using the ``fresh'' randomness in the,
%
hopefully short, %
string $Y$, the function $\ext$ extracts the randomness from $X$ to a nearly perfect form, namely, to a distribution that is close to uniform. We refer to $Y$ as the \emph{seed} of the extractor.

The notion of seeded extractors can be strengthened in different ways. Three such strengthenings that emerged from the study of seeded extractors are strong seeded extractors, non-malleable extractors~\cite{DW09}, and two-source extractors~\cite{CG88}. The latter is the oldest notion, in fact, two-source extractors predate the explicit definition of seeded extractors. Nevertheless, such extractors proved to be the most challenging to construct. In a span of about a decade, strong seeded extractors with nearly optimal parameters were constructed using sophisticated algebraic and combinatorial ideas (see, \eg, \cite{Trevisan01,LRVW03,GUV09,DKSS13,TSU12} as well as~\cite{Shaltiel11} and~\cite[Chapter 6]{Vadhan12}). Non-malleable extractors were introduced more recently, and despite their syntactic resemblance to strong seeded extractors (see \expref{Section}{sec:strong and nm} below for the formal definitions), their constructions required completely different techniques. Furthermore, it was the insight regarding the connection between non-malleable extractors and the seemingly unrelated two-source extractors that enabled the breakthrough
work of Chattopadhyay and Zuckerman~\cite{CZ19} who constructed two-source extractors for
%
polylogarithmic  %
min-entropy. 
%
%
\subsubsection*{A brief and informal summary of our contribution}
%

In this %
article  %
we introduce a new, very natural, variant of seeded extractors we dub \emph{seed-protecting extractors}. Informally, a seeded extractor is seed protecting 
against a class %
$\ccc$ of functions,
mapping seeds to seeds, if the seed $Y$ remains close to uniform even after observing the output $\ext(X,\adv(Y))$ for every choice of $\adv \in \ccc$ (and, of course, a source $X$ with sufficient min-entropy). See \expref{Definition}{def:sp intro} below for the formal definition. We establish what we believe to be surprising and insightful relations, in fact, \emph{equivalences} between seed-protecting extractors and strong extractors, non-malleable extractors against permutations, and two-source extractors, where each case is classified by a suitable class $\ccc$. 
 
 This fresh point of view on non-malleable extractors suggests, in particular, a novel approach for constructing such extractors as, indeed, the focus shifts from breaking output correlations to protecting the seed.  
%
 We first recall the definitions of strong and non-malleable extractors (\expref{Section}{sec:strong and nm}). Then, in \expref{Section}{sec:snsp}, we give the definition of seed-protecting extractors and present out results.

%
%

%

\subsection{Strong seeded extractors and non-malleable extractors}\label{sec:strong and nm}

%
%
\subsubsection{Strong seeded extractors}
%
A $(k,\eps)$-seeded extractor $\ext$ is called \emph{strong} if the output distribution $\ext(X,Y)$ is close to uniform even given the seed $Y$ used for the extraction. This can be expressed by writing
$
(\ext(X,Y),Y) \approx_\eps (U_m,Y).
$
Seeded extractors and, more so, their strong counterparts have found many applications. As mentioned, in a beautiful and deep line of work, efficiently computable strong seeded extractors were constructed for any min-entropy $k$, having seed-length $O(\log\frac{n}{\eps})$ (see \cite{GUV09,DKSS13,TSU12} and references therein). Furthermore,  connections between strong seeded extractors and other objects of study such as list decodable codes, samplers, and expander graphs were found and enabled many applications.%

%
%
\subsubsection{Non-malleable extractors}
%
A \emph{non-malleable extractor} is a strong seeded extractor that has the following additional property. The output of the extractor remains close to uniform even after observing the output of the extractor on any altered seed. Formally, let $\adv \colon \{0,1\}^d \to \{0,1\}^d$ be an arbitrary function with no fixed points, that is, $\adv(y) \neq y$ for all $y$. The reader should think of $\adv$ as an adversarially chosen way of altering the seed. A function $\ext \colon \{0,1\}^n \times \{0,1\}^d \to \{0,1\}^m$ is a \emph{$(k,\eps)$-non-malleable extractor} if for every $(n,k)$-source $X$ and $\adv$ as above,
\begin{equation}\label{eq:nmext}
(\ext(X,Y),\ext(X,\adv(Y)),Y) \approx_\eps (U_m, \ext(X,\adv(Y)),Y)\,,
\end{equation}
where again $Y$ is uniform over $\{0,1\}^d$ and is independent of $X$. %



Non-malleable extractors were introduced by Dodis and Wichs~\cite{DW09}. The original motivation for studying such extractors %
was  %
for the classic problem of devising privacy amplification protocols against active adversaries. Indeed, strong seeded extractors yield a solution to the passive adversary variant. As we discuss later on, non-malleable extractors proved key for the construction of good two-source extractors. More precisely, one requires a certain generalization obtained by considering more than one adversarial function~\cite{CRS12}. Let $t \ge 1$ be an integer. The function $\ext$ above is called a $(k,\eps)$ $t$-non-malleable extractor if for every $t$-tuple of functions $\adv_1,\ldots,\adv_t \colon \{0,1\}^d \to \{0,1\}^d$ with no fixed points, it holds that
\[
(\ext(X,Y),\{\ext(X,\adv_i(Y))\}_{i=1}^{t},Y) \approx_\eps (U_m, \{\ext(X,\adv_i(Y))\}_{i=1}^{t},Y)\,.
\]

In a %
large body of work,
non-malleable extractors were constructed (see \cite{DLWTZ11,CRS12,Li15,Cohen16b,CGL16,Cohen16d,CS16,CL16,Cohen16c,Li17} and references therein). The state-of-the-art construction of $(k,\eps)$-non-malleable extractors~\cite{Li19} has seed length $d = O(\log{n})+ O(\log\frac1{\eps}) \cdot 2^{O(a \cdot (\log\log\frac1{\eps})^{1/a})}$ for min-entropy as low as $k = O(\log\log{n} + a \log\frac1{\eps})$ for every choice of $a \ge 2$. All of these constructions generalize to $t$-non-malleable extractors. Alternatively, a black-box reduction from $t$-non-malleable extractors to non-malleable extractors~\cite{Cohen16d} can be invoked to give explicit $t$-non-malleable extractors with seed length $\poly(t) \cdot d$.

%
%
\subsubsection{Non-malleable extractors against permutations}
%
In this %
article  %
we initiate the study of non-malleable extractors against permutations. These are functions that are ought to satisfy \expeqref{Equation}{eq:nmext} only for  $A$ a permutation with no fixed points. While existing applications of non-malleable extractors (for the construction of %
two-source   %
extractors and for the design of privacy amplification protocols) consider a general adversarial function $A$ with no fixed points, we believe that the new notion of non-malleable extractors against permutations is natural and interesting both in its own right as well for as a step towards constructing full-fledged non-malleable extractors.
It is interesting to note that non-malleable \emph{two-source} extractors against permutations were found useful in
%
independent work  %
by Goldreich and Wigderson~\cite{GW21}, where they
were studied in connection to robustly self-ordered graphs.

\subsection{Strong, non-malleable, and seed-protecting extractors}\label{sec:snsp}
As mentioned, in this %
paper  %
we introduce a new type of randomness extractors which we call \emph{seed-protecting extractors}. To give the formal definition, for an integer $d$, let $\advfam_d$ be the set of all functions from $\{0,1\}^d$ to $\{0,1\}^d$. When $d$ is clear from context, we simply write $\advfam$.%

\begin{definition}[seed-protecting extractors]\label{def:sp intro}
	Let $\ext \colon \{0,1\}^n \times \{0,1\}^d \to \{0,1\}^m$ be a $(k,\eps)$-seeded extractor. Let $\ccc \subseteq \advfam_d$. We say that $\ext$ is \emph{seed protecting against $\ccc$} if for every $\adv \in \ccc$ and every $(n,k)$-source $X$ it holds that
	\begin{equation}\label{eq:intro 123}
	\left(Y, \ext(X,\adv(Y))\right) \approx_\eps \left(U_d, \ext(X,\adv(Y))\right)\,,
	\end{equation}
	where $Y$ is uniformly distributed and is independent of $X$.
\end{definition}



Both non-malleable extractors and seed-protecting extractors have a certain resilience property against tampering %
with  %
the seed. While non-malleable extractors focus on the \emph{output}, seed-protecting extractors are concerned about the \emph{seed}. This shift of focus induces inherent differences. Indeed, while fixed points trivially rule out the possibility of non-malleability (hence, functions with fixed points are excluded by definition), fixed points turn out to be a non-issue for seed-protecting extractors.  Indeed, consider the extreme case -- the identity function. Clearly, non-malleability cannot be achieved against this function. However, note that to be seed protecting against the identity function precisely means to be a strong seeded extractor. In fact, the first observation we make in this preliminary discussion is that strong extractors are equivalent to seed-protecting extractors against
%
a class of permutations %
which we denote by $\perfam \subseteq \mathcal{A}_d$. %

%
\begin{claim}[strong $\iff$ seed protecting against $\perfam$]\label{c:intro strong sp}
Let $\ext \colon \{0,1\}^n \times \{0,1\}^d \to \{0,1\}^m$ be a $(k,\eps)$-seeded extractor. Then,
	\begin{enumerate}
		\item  If $\ext$ is a $(k,\eps)$-seed-protecting extractor against $\perfam$ then $\ext$ is a $(k,2\eps)$-strong seeded extractor.
		\item  If $\ext$ is a $(k,\eps)$-strong seeded extractor then $\ext$ is a $(k,2\eps)$-seed-protecting extractor against $\perfam$.
	\end{enumerate}
\end{claim}
\begin{proof}
  For the first item, a seed-protecting extractor against $\perfam$ can be seen to be a strong %
seeded extractor   %
by taking $\adv$ to be the identity function. Indeed, with this choice, since $\ext$ is
a 
%
seeded extractor   %
the right hand side of \expeqref{Equation}{eq:intro 123} is $\eps$-close to $U_{d+m}$, and so $(Y,\ext(X,Y)) \approx_{2\eps} (Y,U_m)$. 	For the other direction, as $\ext$ is $(k,\eps)$-strong we have that $(Y,\ext(X,Y)) \approx_\eps (Y,U_m)$. Thus, for any $\adv \in \perfam_d$, it holds that 
	$$(\adv(Y),\ext(X,\adv(Y))) \approx_\eps (\adv(Y),U_m)\,.$$
	By the data processing inequality, one can apply any function to the first component $\adv(Y)$ of both sides and maintain the $\eps$-closeness. In particular, by applying $\adv^{-1}$ we get 
	\begin{equation}\label{eq:12221121211}
	(Y,\ext(X,\adv(Y))) \approx_\eps (Y,U_m)\,.
	\end{equation}
	Now, $(Y,U_m)$ has the same distribution as $U_{d+m}$, and $\ext(X,\adv(Y)) \approx_\eps U_m$. This, together with \expeqref{Equation}{eq:12221121211}, implies that $(Y,\ext(X,\adv(Y))) \approx_{2\eps} (U_d,\ext(X,\adv(Y)))$, completing the proof.
\end{proof}

Going back to non-malleable extractors, by the discussion above, it is not a priori clear whether non-malleability is in any way related to seed protection. Nonetheless, one of the results of %
this paper %
is an \emph{equivalence} between the property of non-malleability and seed protection, at least when focusing on permutation adversaries. By saying that $\ext$ is a non-malleable extractor against $\ccc \subseteq \advfam$, we mean that \expeqref{Equation}{eq:nmext} holds for every $\adv \in \ccc$ (but not necessarily for other functions) that, in addition, has no fixed points. 

For stating our result, we generalize seed protection to several adversarial functions. This should be done with some care. Indeed, naively, one's first suggestion might require that for every two functions $\adv_1,\adv_2 \in \ccc$, it holds that
\[
\left(Y, \ext(X,\adv_1(Y)),\ext(X,\adv_2(Y))\right) \approx_\eps \left(U_d, \ext(X,\adv_1(Y)),\ext(X,\adv_2(Y))\right)\,.
\]
This definition, we observe, is moot. Indeed, consider two functions $A_1, A_2 \colon \{0,1\}^d \to \{0,1\}^d$ that according to the first bit of the seed, $Y_1$, decide whether to ``behave'' exactly the same or very differently. More concretely, sample two permutations $\pi_0,\pi_1$ on $\{0,1\}^d$ at random, and define $A_1(y) = \pi_0(y)$ and $A_2(y) = \pi_{y_1}(y)$.
As $\pi_0,\pi_1$ were chosen at random, and thus are usually disagree, by observing $\ext(X,\adv_1(Y))$ and $\ext(X,\adv_2(Y))$ one can distinguish $Y$, in fact its first bit $Y_1$, from uniform by checking whether both outputs are equal.

As we prove, this %
``collusion,''  %
%
in which two or more  %
adversarial functions attain the same value, is the only obstacle for seed protection against permutations.  %
%
(Already  %
here we stress that there are other obstacles when considering functions other than permutations, as we discuss
in \expref{Section}{sec:intro two source sp}.) %
%
Given $\adv_1,\ldots,\adv_t \colon \{0,1\}^d \rightarrow \{0,1\}^d$, we say $\adv_1,\ldots,\adv_t$ are
\emph{non-colluding} if for every $y \in \{0,1\}^d$, all the evaluations $\adv_1(y),\ldots,\adv_{t}(y)$ are distinct. We denote by $\nc^t \subseteq \advfam^t_d$ the set of $t$-tuples of functions that are non-colluding.
%
%
%
%
With hindsight, we give the following generalization of seed-protecting extractors to several adversarial functions.

\begin{definition}
Let $\ext \colon \{0,1\}^n \times \{0,1\}^d \to \{0,1\}^m$ be a $(k,\eps)$-seeded extractor. Let $\ccc \subseteq \advfam$. We say that $\ext$ is \emph{$t$-seed protecting against $\ccc$}
if for every $(n,k)$-source $X$ and $(\adv_1,\ldots,\adv_t) \in \nc^t \cap \ccc^t$ it holds that
\[
\left(Y, \{\ext(X,\adv_i(Y))\}_{i=1}^{t}\right) \approx_\eps \left(U_d, \{\ext(X,\adv_i(Y))\}_{i=1}^{t}\right)\,.
\]
We also express this by saying that $\ext$ is seed protecting against $\ccc^t$.
\end{definition}

The following lemma gives the equivalence's easier direction, showing that non-malleable extractors against permutations (with no fixed points) are seed protecting against (non-colluding) permutations. We refer the reader to \expref{Lemma}{lem:nm-to-sp} for a more general statement.

\begin{lemma}[non-malleable $\Longrightarrow$ seed protecting]\label{lem:intro nm is sp}
	Let $t \ge 1$ and assume $\ext$ is a $(k,\eps)$-non-malleable extractor
	against $\Pi^t$. Then, $\ext$ is a $(k,4t\eps)$-seed-protecting extractor against $\Pi^{t+1}$.
\end{lemma}

As a %
warm-up,  %
in \expref{Section}{sec:overview}, we prove \expref{Lemma}{lem:intro nm is sp} for $t=1$ as well as its converse which is indeed more surprising and difficult to prove (see \expref{Theorem}{thm:warmup}). 

Towards stating the other direction, 
for $\Delta \ge 0$, define $\mathcal{F}_{\Delta}$ to be the subset of $\advfam_d$ containing all functions $A \colon \{0,1\}^d \to \{0,1\}^d$ such that 
$\me(A(U_d)) \ge d-\Delta$. %
%
%
For example, note that $\mathcal{F}_0 = \perfam$. We
give a reduction from non-malleable extractors against $\mathcal{F}_\Delta$ to seed-protecting extractors against $\mathcal{F}_\Delta$. 

\begin{theorem}[seed protecting $\Longrightarrow$ non-malleable]\label{thm:last one}
	Let $t \ge 1$ be an integer, and $\Delta \ge \max(1,\log{t})$.
	Let $\ext\colon\B^{n}\times\B^{d}\rightarrow\B^m$ be a $(k,\eps)$-seed-protecting extractor against $\mathcal{F}^{t+1}_{\Delta} \cap \nc^{t+1}$. Assume further that $d = \Omega(\log{t})$. Then, $\ext$ is $(k',\eps')$-non-malleable against $\mathcal{F}_{\Delta}^t$ with $k' = k+mt+\log\frac1{\eps}$ and $\eps' = O(\eps^{1/3})$.
\end{theorem}

We prove \expref{Theoren}{thm:last one} in \expref{Section}{sec:the delta proof}. 
We note that for $\Delta = 0$ and $t = 1$, seed protection against
permutations is enough, and one does not need to devise an extractor against $\mathcal{F}_1$.
Thus, for $\Delta = 0$ and $t = 1$ there is a strong equivalence, which we prove 
as a warm-up, in \expref{Theorem}{thm:warmup}.


%

%

 %


%

%

\subsection{Two-source extractors as seed-protecting extractors}\label{sec:intro two source sp}

Discussing seed protection against permutations sufficed for characterizing both strong and non-malleable extractors against permutations. 
But, what about other adversarial functions? Is it the case that seed protection is achievable against any single function?
%
(Of course, collusion is irrelevant in such a setting.) %
The quick answer is %
``no.''   %
Surprisingly, our next result is a characterization we obtain for two-source extractors as $1$-seed-protecting extractors %
\ie,  %
seed-protecting extractors with a single adversarial function) against a suitable family. In particular, known impossibility results on two-source extractors translate to impossibility results on 1-seed-protecting extractors. %

%

Before recalling the formal definition of two-source extractors and describing this family, we believe it is instructive to first consider an extreme case and ask whether one can seed-protect against an adversarial function $\adv$ that, unlike a permutation, is allowed to ``focus'' on seeds of its choice.  The ultimate case is where $\adv$ has range of size one.  %
%
%
However, in such case, $\adv(y)$ gives no information about the seed $y$, and so seed %
protection  %
trivially follows. What about
a  %
range of size two? We have the following easy claim that establishes
the  %
%
impossibility of seed protection against such functions. Let $\advpair \subseteq \mathcal{A}_d$ be the set of all functions $\adv \colon \{0,1\}^d \to \{0,1\}^d$ with range of size precisely two.
%

\begin{claim}\label{c:lower bound image two intro}
	Let $\ext \colon \{0,1\}^n \times \{0,1\}^d \to \{0,1\}$ be a $(k,\eps)$-seeded extractor. Then, for $k \le n-1$ and $\eps < \frac1{6}$, $\ext$ is not $(k,\eps)$-seed protecting against $\advpair$.
\end{claim}
\begin{proof}
	Fix an arbitrary $w \in \{0,1\}^d$ and assume without loss of generality that $\Pr[\ext(U_n,w) = 0] \ge \frac1{2}$. Define $X$ to be the random variable that is uniformly distributed over all $x \in \{0,1\}^n$ such that $\ext(x,w ) = 0$. Since $\ext$ is a 
%
$(k,\eps)$-seeded extractor   %
for min-entropy $k$ and as $\me(X) \ge n-1 \ge k$, we have that $\ext(X,Y) \approx_\eps U_1$. Thus,
	$$
	\Pr[\ext(X,Y) = 1] \ge \frac1{2} - \eps\,.
	$$
	By an averaging argument, there exists $z \in \{0,1\}^d$ such that
	$$
	\Pr[\ext(X,z) = 1] \ge \frac1{2} - \eps\,.
	$$
	Note that $w \neq z$. Define the function $\adv : \{0,1\}^d \to \{0,1\}^d$ by
	\begin{align*}
	\adv(y) = \begin{cases}
	w & y_1 = 0,\\
	z & y_1 = 1,
	\end{cases}
	\end{align*}
	and note that $\adv \in \advpair$. Denote $Z = \ext(X,Y)$ and $Z' = \ext(X,A(Y))$.	We turn to show that
	$$
	\gamma = \SD\left(\left(Y,Z'\right) , \left(U_d, Z'\right)\right) \ge \frac1{4} - \frac{\eps}{2}\,.
	$$
	To see this, note that
	\begin{equation}\label{eq:12331}
	\gamma \ge \SD\left(\left(Y_1, Z'\right) , \left(U_1,Z'\right)\right)
	\ge \Pr[Z' = Y_1] - \Pr[Z' = U_1]\,.
	\end{equation}
	%
	We have that
	\begin{align*}
	\Pr[Z' = Y_1] &=
	\frac1{2} \Pr[Z' = Y_1 \mid Y_1 = 0] + \frac1{2} \Pr[Z' = Y_1 \mid Y_1 = 1] \\
	&= \frac1{2} \left(\Pr[\ext(X,w) = 0] + \Pr[\ext(X,z) = 1] \right) \\
	&\ge \frac1{2} \left(1 + \frac1{2} - \eps\right) \\
	&= \frac{3-2\eps}{4}\,.
	\end{align*}
	On the other hand, $\Pr[Z' = U_1] = \frac1{2}$, and so by Equation~\eqref{eq:12331}\,, $$
	\gamma \ge \frac{3-2\eps}{4} - \frac1{2} \ge \frac1{4} - \frac{\eps}{2}\,.
	$$
	Thus, as we assume $\eps < \frac1{6}$, we get that
	$\gamma > \frac1{6}$ which concludes the proof.
\end{proof}

We next recall the definition of a two-source extractor or, more generally, of
an %
unbalanced two-source %
%
extractor,  %
and then present our characterization of two-source extractors as seed-protecting extractors. A \emph{$(k_1,k_2,\eps)$-two-source extractor} is a function $\ext \colon \{0,1\}^{n_1} \times \{0,1\}^{n_2} \to \{0,1\}^m$ such that for every $(n_1,k_1)$-source $X$ and an independent $(n_2,k_2)$-source $Y$ it holds that $\ext(X,Y) \approx_\eps U_m$. The existence of a two-source extractor for min-entropies $k_1 = \log{n_2}+ O(\log\frac1{\eps})$ and $k_2 = \log{n_1} + O(\log\frac1{\eps})$ with $m=k_1+k_2-O(\log\frac1{\eps})$ output bits was established %
%
in~\cite{CG88}  %
but the problem of explicitly constructing a $(k,k,\eps)$-two-source extractor with $n_1 = n_2 = n$ even for min-entropy as high as $k = 0.49n$ remained open for three decades \cite{CG88,Bourgain05,Raz05}.

Over the last few years, there has been remarkable progress on this problem. In particular, in a breakthrough %
paper %
Chattopadhyay and Zuckerman~\cite{CZ19}  obtained a $(k,k,\frac1{n})$-two-source extractor $\ext \colon \{0,1\}^n \times \{0,1\}^n \to \{0,1\}$ for min-entropy as low as $k = \poly(\log{n})$. %
Subsequent  %
work
\cite{BADTS17,Cohen17,Li19} improved the entropy requirement even further to $k = \widetilde{O}(\log{n})$. Constructing two-source extractors for min-entropy $O(\log{n})$ (or, more ambitiously, $\log(n)+O(1)$) is highly motivated by the problem of constructing explicit Ramsey graphs \cite{BKSSW10,BRSW12,Cohen16,Cohen17}. A second important open problem is constructing two-source extractors with low error. Current techniques do not yield explicit two-source extractors when 
%
%
$\eps = 1/n^{\omega(1)}$. 
%



Our main result here is proving an equivalence of two-source extractors and seed-protecting extractors for the class $\mathcal{F}_{\Delta}$.  %
%
%


%
\begin{theorem}[two-source extractors as seed-protecting extractors]  %
\label{thm:two source sp intro}
	Let $\ext \colon \{0,1\}^n \times \{0,1\}^d \to \{0,1\}^m$. 
	\begin{enumerate}
		\item If $\ext$ is $(k_1,\eps)$-seed-protecting extractor against $\mathcal{F}_{d-k_2}$ then $\ext$ is a $(k_1,k_2,3\eps)$-two-source extractor.
		\item If $\ext$ is a $(k_1,k_2,\eps)$-two-source extractor which is strong in the second source\footnote{Note that this %
strength %
requirement has the undesired effect of breaking the equivalence. However, one can always assume that a two-source extractor is strong in each of its sources provided one is willing to increase the error by a multiplicative factor of $2^{O(m)}$. See also a remark in the footnote of~\cite{Raz05} just prior to Definition 1.3.} then $\ext $ is $(k_1,2\eps)$-seed protecting against $\mathcal{F}_{d-k_2}$.
	\end{enumerate}
\end{theorem}

\expref{Theorem}{thm:two source sp intro} is proven in \expref{Section}{sec:two source} (see \expref{Lemma}{lem: sp is ts} and \expref{Lemma}{lem:two-source} for the proof of each direction. %
By invoking a known lower bound result on the amount of min-entropy required for two-source extractors, \expref{Theorem}{thm:two source sp intro} in particular implies that seed protection cannot be achieved against adversaries that are allowed to have small range. \expref{Claim}{c:lower bound image two intro} gives a direct proof of that for
%
%
the  %
extreme case of range size two. We find it insightful that the
%
natural notion of seed protection gives a characterization of the three most well-studied types of randomness extractors: strong, two-source, and non-malleable (albeit against adversaries with large range).



\subsection{Open problems}

For which other classes $\mathcal{C}$ do seed-protecting extractors against $\mathcal{C}$ exist?
Can we get full-fledged non-malleability (\ie,
against general adversaries with no fixed points) from seed-protecting extractors against these prospective %
%
classes $\mathcal C\,$?
More generally, extending the connection between seed protecting and non-malleable extractors is an intriguing open question left for future research.

%
%
%
%
%

For $\mathcal{C} = \Pi$, \ie, for permutation adversaries, we proved that indeed seed protection implies 
non-malleability. An interesting open problem is whether non-malleable extractors for permutations (or even involutions) imply full-fledged non-malleability.
A second interesting open problem is whether non-malleability against permutations (or involutions) implies that for each source $X$ there exists a small set $B_X \subset \{0,1\}^d$ of
``bad'' seeds such that $(\ext(X,y_1),\ext(X,y_2)) \approx U$ for every two distinct $y_1,y_2 \in \{0,1\}^d \setminus B_X$. Note that this 
weaker notion of non-malleability already suffices for Chattopadhyay and Zuckerman's construction of %
two-source  %
extractors~\cite{CZ19}.

Also, note that \expref{Theorem}{thm:last one} incurs $tm$ entropy loss. Recall that non-malleable extractors must satisfy $k \ge (t+1)m$, as the $t+1$ outputs must be independent for a proper choice of adversaries. On the other hand, intuitively, this requirement is unnecessary for seed-protecting extractors, and it seems that the only requirement should be $k \ge m$ (we ignore the additive error dependence for simplicity).  Having said that, we do not know how to formalize the intuition above regarding the entropy loss of seed-protecting extractors. Indeed, in \expref{Section}{sec:nexp} we prove the existence of seed-protecting extractors via a probabilistic argument, and our proof technique requires $k \ge tm$. We leave the question of understanding the entropy loss of seed-protecting extractors to future research.

Lastly, it would be interesting to investigate the connection between seed-protecting extractors to other types of extractors. Concretely, can \emph{non-malleable two-source} extractors be characterized by seed-protecting extractors?

%



\section{Warm-up}\label{sec:overview}

As a warm-up, in this section we give a proof sketch for the equivalence between non-malleable extractors and seed-protecting extractors for permutations. For simplicity we focus on the case $t = 1$. The case $t > 1$ follows by similar ideas but is somewhat more involved. 

\begin{theorem}[$\perfam$ non-malleability $\iff$ $\perfam^2$ seed protection]\label{thm:warmup}
Let $\ext \colon \{0,1\}^n \times \{0,1\}^d \to \{0,1\}^m$.%
\begin{enumerate}
	\item If $\ext$ is $(k,\eps)$-seed protecting against $\perfam^2$ then $\ext$ is $(k',14\eps^{1/3})$-non-malleable against $\perfam$, where $k' = k+2m+O(\log(1/\eps))$.
	\item If $\ext$ is $(k,\eps)$-non-malleable against $\perfam$ then $\ext$ is $(k,3\eps)$-seed-protecting extractor against $\perfam^2$.
\end{enumerate}
\end{theorem}

\begin{proof}
We start with the first and more difficult item.
%
\subsubsection*{Proof sketch of the first item}
Set $\delta =14\eps^{1/3}$. Assume towards a contradiction that $\ext$ is not non-malleable against $\perfam$. Consider then a source $X\sim \B^n$ and adversarial permutation $\adv \in \Pi$ with no fixed points for which
	\begin{align*}
	\SD \Big( \big( \ext(X,Y), \ext(X,\adv(Y)), Y \big), 
	\big( U_m, \ext(X,\adv(Y)),  Y \big)
	\Big) > \delta\,,
	\end{align*}
 where $Y \sim \B^d$ is uniform and independent of $X$.
	This implies
	\begin{align}\label{eq:not nm i}
	\E_{y \sim Y}\Big[ \SD\big( ( \ext(X,y), \ext(X,\adv(y)) ), 
	( U_m, \ext(X,\adv(y)) )\big) \Big] > \delta\,.
	\end{align}
	For every $y,y_1 \in \{0,1\}^d$ we define the distributions 
	\begin{align*}
	\mathcal{D}_{y,y_1} &= \left( \ext(X,y), \ext(X,y_1)  \right), \\
	\mathcal{I}_{y_1}   &= \left(U_m, \ext(X,y_1) \right)\,.
	\end{align*}
	Define the function $T \colon \{0,1\}^d \times \{0,1\}^d \to [0,1]$ by
	$
	T(y,y_1) = \SD\left( \mathcal{D}_{y,y_1}, \mathcal{I}_{y_1} \right).
	$
	With this notation, \expeqref{Equation}{eq:not nm i} can be written as
	$
	\E\limits_{y \sim Y}{\left[T(y,\adv(y))\right]} > \delta.
	$
	By an averaging argument, there exists a set $H \subseteq \{0,1\}^d$ of size 
	%
	%
	$|H| = \delta/2 \cdot 2^d$ such that
	%
	%
	%
	$T(y,\adv(y)) > \delta/2$ 
	%
	for every $y \in H$.
	
	Based on the fact that $\ext$ is strong, we prove that  	
	there exists a subset $B_1 \subseteq \{0,1\}^d$ of density $2\eps^{1/3}$ such that for every $y \notin B_1$, $
		\E_{y_1 \sim U_{d}}\left[T(y,y_1)\right] \le \eps^{2/3}.$ We remark that this is where one pays $2m+O(\log(1/\eps))$ in the entropy loss. We choose to skip the proof of this fact (see \expref{Claim}{claim:average behavior}).

                A main part of the proof is extending $A|_H$, the restriction of $A$ to $H$, to a new permutation $\widehat{A}$ over $\{0,1\}^d$ such that for almost every $y$ outside of $H$ it holds that $T(y,\widehat{A}(y))$ is small, in particular, bounded by $\eps^{1/3}$. This is done via a greedy algorithm. We arrange the elements of $\{0,1\}^d \setminus (H \cup B_1)$ in some order $y_1,\ldots,y_\ell$. By an averaging argument, for every $y \not\in B_1$, there are at most $\eps^{1/3}$ fraction of seeds $y_1$ for which $T(y,y_1) \ge \eps^{1/3}$. Denote this set by $B(y)$. We proceed iteratively, starting from $i = 1$, and choose an element $z_i \not\in B(y_i)$, also different from $y_i$, that has not been assigned already as %
an element of the range %
of (the partially defined) $\widehat{A}$, 
  %
and set $\widehat{A}(y_i) = z_i$. This can be done for most
%
elements $y_i$.  %
%
%
When $i$ gets very close to $\ell$ we may have to assign the remaining 
%
elements of the range %
in any way, but which will still
%
 guarantee that $\widehat{A}$ is a permutation. At any rate, for simplicity, let us assume that for all elements $y_1,\ldots,y_\ell$ we have that $\widehat{A}(y_i) \not\in B(y_i) \cup \{y_i\}$.
		

Define the random variables $\mathcal{D}_{\widehat{A}} = \mathcal{D}_{Y,\widehat{A}(Y)}$ and 
$\mathcal{I}_{\widehat{A}} = \mathcal{I}_{\widehat{A}(Y)}$,
	where $Y$ is uniformly distributed over $\{0,1\}^d$. Using the fact that $\ext$ is seed protecting, we prove the following.
	
	\begin{claim}\label{claim:d is concentrated i}
		There exists a set $B_2 \subseteq \{0,1\}^d$ of density at most $\sqrt{\eps}$ such that for every $y \in \{0,1\}^d \setminus B_2$,
		$
		\SD(\mathcal{D}_{\widehat{A}}, \mathcal{D}_{y,\widehat{A}(y)}) \le \sqrt{\eps}.
		$
	\end{claim}
	
	\begin{proof}
		%
		Denote 
\[
		\widehat{Z}(y) =\ext(X,\widehat{A}(y))\,.
\]
		With this notation, we have that
		\begin{align*}
		\mathcal{D}_{\widehat{A}} &= ( \ext(X,Y), \widehat{Z}(Y) ),\\
		\mathcal{D}_{y,\widehat{A}(y)} &= ( \ext(X,y), \widehat{Z}(y) )\,.
		\end{align*}
		%
			By construction, $\widehat{A}$ is a permutation with no fixed points (more importantly, it does not collude with the identity function). Observe that as $\ext$ is seed protecting against $\perfam^2 \cap \nc^2$ with error $\eps$, 
		\[
		\E_{y\sim U_d}[\SD(\mathcal{D}_{y,\widehat{A}(y)}, \mathcal{D}_{\widehat{A}})] \leq \eps.
		\]
By Markov's inequality, the set $B_2 \subseteq \{0,1\}^d$ of %
all the $y$  %
satisfying
$
		\SD((\mathcal{D}_{y,\widehat{A}(y)}, \mathcal{D}_{\widehat{A}})) \geq \sqrt{\eps}
		$
		has density at most $\sqrt{\eps}$, as stated.
		%
		
	\end{proof}
	Using that $\ext$ is a strong seeded extractor, one can prove that
	\begin{claim}\label{claim:i is concentrated i}
		There exists a set
		$B_3 \subseteq \{0,1\}^d$ of density at most $\sqrt{\eps}$ such that for every $y \in \{0,1\}^d \setminus B_3$,
		$
		\SD(\mathcal{I}_{\widehat{A}}, \mathcal{I}_{\widehat{A}(y)}) \le \sqrt{\eps}.
		$
	\end{claim}
	
	Write $B = B_1 \cup B_2 \cup B_3$. Recall that 
	%
	%
	$\abs{H} = \delta/2\cdot 2^d = 7\eps^{1/3} \cdot 2^d.$ 
	%
	By the above claims we can bound
	$
	\abs{B} < 7\eps^{1/3} \cdot 2^d,
	$
	and so there exists $y_h \in H \setminus B$.  Take $y_\ell \in \{0,1\}^d \setminus (H \cup B)$. 	
	By \expref{Claim}{claim:d is concentrated i}, since $y_\ell,y_h \notin B_2$,
	$\SD(\mathcal{D}_{y_h,\widehat{A}(y_h)}, \mathcal{D}_{\widehat{A}} ) \le \sqrt{\eps}$ and
	$\SD(\mathcal{D}_{y_\ell,\widehat{A}(y_\ell)}, \mathcal{D}_{\widehat{A}} ) \le \sqrt{\eps}.
	$
	Hence,
	\begin{equation}\label{eq:dldh i}
	\SD\left(\mathcal{D}_{y_\ell,\widehat{A}(y_\ell)}, \mathcal{D}_{y_h,\widehat{A}(y_h)}\right) \le 2\sqrt{\eps}\,.
	\end{equation}
	By \expref{Claim}{claim:i is concentrated i} and since $y_\ell,y_h \notin B_3$,
 	$
	\SD(\mathcal{I}_{\widehat{A}(y_h)}, \mathcal{I}_{\widehat{A}} ) \le \sqrt{\eps}$ and
	$\SD(\mathcal{I}_{\widehat{A}(y_\ell)}, \mathcal{I}_{\widehat{A}} ) \le \sqrt{\eps}.
$
	Thus,
	\begin{equation}\label{eq:ilih i}
	\SD\left(\mathcal{I}_{\widehat{A}(y_\ell)}, \mathcal{I}_{\widehat{A}(y_h)} \right) \le 2\sqrt{\eps}\,.
	\end{equation}
	Now, since $y_\ell \in \{0,1\}^d \setminus (H \cup B_1)$,
	\begin{equation}\label{eq:dlil i}
	\SD\left(\mathcal{D}_{y_\ell,\widehat{A}(y_\ell)}, \mathcal{I}_{\widehat{A}(y_\ell)}\right) \le \eps^{1/3}\,.
	\end{equation}
%
 By \expeqref{Equations}{eq:dldh i}, \eqref{eq:ilih i}, \expeqref{and}{eq:dlil i}, %
and the triangle inequality,
	$
	\SD(\mathcal{D}_{y_h,\widehat{A}(y_h)}, \mathcal{I}_{\widehat{A}(y_h)}) \le 5 \eps^{1/3}.
	$
	However, $y_h \in H$ and so
	\begin{align*}
	\SD\left( \mathcal{D}_{y_h,\widehat{A}(y_h)}, \mathcal{I}_{\widehat{A}(y_h)} \right) = T\left(y_h,\widehat{A}(y_h)\right) 
	= T(y_h, \adv(y_h)) > \frac{\delta}{2}\,,
	\end{align*}
	contradicting our choice of $\delta$. This proves the first item.
	
%
\subsubsection*{Proof sketch of the second item}
	
Moving on to the second item, let $A_1,A_2 \in \Pi$ be %
non-colluding.
Let $X$ be an $(n,k)$-source, and let $Y \sim \B^{d}$ be 
	a uniform random variable, independent of $X$. Denote	
	\[
	Z(y) = ( \ext(X,A_{1}(y)),\ext(X,A_{2}(y)) ).
	\]
	As $A_{1}^{-1}$ is a permutation, $A_1^{-1}(Y)$ distributes the same as $Y$ does. Thus,
	\[
	\SD((U_d,Z(Y)),(Y,Z(Y))) = 
	\SD((U_d,Z(A_{1}^{-1}(Y))),({A}_{1}^{-1}(Y),{Z}({A}_{1}^{-1}(Y)))).
	\]
	By the data-processing inequality,\footnote{By this we refer to the fact that for any random variables
	$X,Y \sim \Omega_1$, and every $f \colon \Omega_1 \rightarrow \Omega_2$, possibly randomized, $\SD(f(X),f(Y)) \le \SD(X,Y)$. Equality is attained when $f$ is injective.}
	we can apply ${A}_1$ on the prefix of both random variables without increasing
	the statistical distance; The above equation then becomes
	\begin{equation}\label{eq:sec423 i}
	\SD((U_d,{Z}(Y)),(Y,{Z}(Y))) = 
	\SD((U_d,{Z}({A}_{1}^{-1}(Y))),(Y,{Z}({A}_{1}^{-1}(Y))))\,.
	\end{equation}
	Define $P = {A}_{2} \circ {A}_{1}^{-1}$. Then,
	\begin{align*}
	{Z}({A}_{1}^{-1}(Y)) = 
	\left( \ext\left( X,Y \right),\ext( X,P(Y)) \right)\,.
	\end{align*}
	Note that $P$ is a permutation. Furthermore, observe that as $({A}_1,{A}_{2}) \in \nc^{2}$,  $P$ has no fixed points. Indeed, assuming
	towards a contradiction that $P(y)=y$ for some $y \in \B^{d}$, we get
	that ${A}_{1}^{-1}(y) = {A}^{-1}_{2}(y)$, which is impossible.
	By the non-malleability of $\ext$, we know that
	\[
	\SD((Y,{Z}({A}_1^{-1}(Y))),(Y,U_m,\ext(X,P(Y)))) \le \eps\,,
	\]
    Thus, by the triangle inequality, we get
	\begin{align}
	\SD((U_d,{Z}({A}_{1}^{-1}(Y))),(Y,{Z}({A}_{1}^{-1}(Y)))) &\le 
	\SD((U_d,{Z}({A}_{1}^{-1}(Y))),(U_d,U_m,\ext(X,P(Y)))) + \nonumber \\
	& ~~~~\SD((U_d,U_m,\ext(X,P(Y))),(Y,U_m,\ext(X,P(Y)))) + \nonumber \\
	& ~~~~\SD((Y,U_m,\ext(X,P(Y))),(Y,{Z}({A}_{1}^{-1}(Y)))) \nonumber \\
	&\le 3\eps\,.
	\end{align}
	The proof then follows by \expeqref{Equation}{eq:sec423 i}.
\end{proof}

%
\section{Seed-protecting  %
extractors for high entropy adversaries are non-malleable}\label{sec:the delta proof}

In this section we prove \expref{Theorem}{thm:last one}, showing that a seed-protecting extractor with adversarial entropy parameter $\Delta$
is non-malleable against roughly the same adversarial function class.



\begin{definition}[non-colluding functions]
	We say that a tuple $(\adv_1,\ldots,\adv_t) \in \advfam_d^t$ 
	\emph{do not collude}, if for every $y \in \B^d$, $\adv_1(y),\ldots,\adv_{t}(y)$ are pairwise distinct. 
	When $d$ is clear from context, we denote by $\nc^t$ the set of $t$-tuples of non-colluding
	functions from $\B^d$ to $\B^d$. %
 	%
\end{definition}



We prove the following slight restatement of \expref{Theorem}{thm:last one}. Henceforth, we denote by
$\nofp \subseteq \advfam_d$ the class of functions with no fixed points. Recall that $\mathcal{F}_{\Delta}$ is the set of functions $A \in \advfam_d$ for which $\me(A(U_d)) \ge d-\Delta$.

\begin{theorem}\label{t: sp to nm}
Let $t \ge 1$ be an integer, and $\Delta \ge 0$.
Let $\ext\colon\B^{n}\times\B^{d}\rightarrow\B^m$ be a $(k,\eps)$-seeded extractor with $d \ge 2\log{t} + \log\frac1{\eps} + 2$ and $\eps \leq 10^{-4}$. Assume that $\ext$ is seed protecting against $\advfampart^{t+1}_{\max(\Delta, \log t,1)}$. Then, $\ext$ is non-malleable against $\left(\advfampart_{\Delta} \cap \nofp\right)^t \cap \nc^t$  for min-entropy $k+mt+\log\frac1{\eps}$ and error guarantee $14 \eps^{1/3}$.
\end{theorem}

Note that the non-malleability guaranteed by \expref{Theorem}{t: sp to nm} is with respect to $\left(\advfampart_{\Delta} \cap \nofp\right)^t \cap \nc^t$. That is, non-colluding is still required. In \expref{Theorem}{sec:no nc} we show how to get rid of this requirement in a ``black-box'' fashion.

\begin{proof}[Proof of \expref{Theorem}{t: sp to nm}.]
Assume towards a contradiction that $\ext$ is not non-malleable against $\left(\advfampart_{\Delta} \cap \nofp\right)^t$ for min-entropy $k' = k + mt + \log(1/\eps)$. Then, there exists an $(n,k')$ source $X$ and functions $\adv_1,\ldots,\adv_t \in \advfampart_{\Delta} \cap \nofp$ with $(\adv_1,\ldots,\adv_t) \in \nc^t$, such that
\begin{align*}
\SD \Big( \big( \ext(X,Y), &\ext(X,\adv_1(Y)),\ldots, \ext(X,\adv_t(Y)), Y \big), \\
           \big( U_m, &\ext(X,\adv_1(Y)),\ldots, \ext(X,\adv_t(Y)),  Y \big)
    \Big) > \delta\,,
\end{align*}
for $\delta = 14 \eps^{1/3}$, where $Y \sim \B^d$ is uniform and independent of $X$.
That is,
\begin{align}\label{eq:not nm}
\E_{y \sim Y}\Big[ \SD\big( ( \ext(X,y), &\ext(X,\adv_1(y)), \ldots, \ext(X,\adv_t(y)) ), \notag \\
           ( U_m, &\ext(X,\adv_1(y)), \ldots, \ext(X,\adv_t(y)) )\big) \Big] > \delta\,.
\end{align}
For $y,y_1,\ldots,y_t \in \{0,1\}^d$ define the distributions 
\begin{align*}
\mathcal{D}_{y,y_1,\ldots,y_t} &= \left( \ext(X,y), \ext(X,y_1),\ldots,\ext(X,y_t)  \right), \\
\mathcal{I}_{y_1,\ldots,y_t}   &= \left(U_m, \ext(X,y_1),\ldots,\ext(X,y_t) \right)\,.
\end{align*}
Define the function $T \colon \left(\{0,1\}^d\right)^{t+1} \to [0,1]$ as follows. For $y,y_1,\ldots,y_t \in \{0,1\}^d$,
\[
T(y,y_1,\ldots,y_t) = \SD\left( \mathcal{D}_{y,y_1,\ldots,y_t}, \mathcal{I}_{y_1,\ldots,y_t} \right)\,.
\]
With this notation, \expeqref{Equation}{eq:not nm} can be written as
$
\E\limits_{y \sim Y}{\left[T(y,\adv_1(y), \ldots, \adv_t(y))\right]} > \delta.
$
By an averaging argument, there exists a set $H \subseteq \{0,1\}^d$ of size $|H| = (\delta/2) \cdot 2^d$ such that for every $y \in H$,
\[
T(y,\adv_1(y), \ldots, \adv_t(y)) > \frac{\delta}{2}\,.
\]

 
Let $Y_1,\ldots, Y_t$ be independent random variables, that are jointly 
independent of $X$, and each $Y_i$ is uniformly distributed over 
$\{0,1\}^d$. Denote 
\[
Z = \ext(X,Y_1)\circ\cdots\circ\ext(X,Y_t)\,.
\] For 
$z = (z_1,\ldots , z_t)  \in (\B^m)^t$ define the random variable $X_{z} = X \mid \{Z=z\}.$
\begin{claim}\label{lem: cond entropy most cases}
There exists a set $B' \subseteq (\{0,1\}^m)^t$ such that: 
\begin{enumerate}
	\item\label{it1} $\Pr[Z \in B'] \le \eps$, and,
	\item\label{it2} For every $z \notin B'$, $\minent{X_{z}}  \ge \minent{X} - tm - \log(1/\eps) \ge k$.
\end{enumerate}
\end{claim}

\begin{proof}
%
	
	Fix $z\in (\{0,1\}^m)^t$ and observe that  
	\begin{equation}\label{eq: condition min entropy loss delta}
		\minent{ X_{z}} \geq \minent{X} - \log\left(\frac{1}{\Pr[Z = z]}\right)\,.
	\end{equation}
	Define
	$
	B' = \set{z : \Pr[Z = z] \le 2^{-mt}  \eps}.
	$
	By \expeqref{Equation}{eq: condition min entropy loss delta}, for every $z \notin B'$, \expeqref{item}{it2} holds.
 	As $B' \subseteq \{0,1\}^{mt}$,
\begin{align*}
	\Pr[Z\in B'] &= \sum_{z \in B'}\Pr[Z = z] 
	\leq |B'| \cdot 2^{-mt}  \eps \le \eps\,,
\end{align*}
and so \expeqref{item}{it1} follows as well.
\end{proof}

\begin{claim}\label{claim:average behavior}
There exists a subset $B_1 \subseteq \{0,1\}^d$ of size $|B_1| \le 2\eps^{1/3} \cdot 2^d$ such that for every $y \notin B_1$,
\[
\E_{(y_1,\ldots,y_t) \sim U_{td}}\left[T(y,y_1,\ldots,y_t)\right] \le \eps^{2/3}\,.
\]
\end{claim}

\begin{proof}
Recall the definition of $X_{z}$ and $B'$ from \expref{Lemma}{lem: cond entropy most cases}. 
Fix $z \not\in B'$. Let $Y$ be uniform over $\{0,1\}^d$ and independent of $(X,Y_1,\ldots, Y_t)$. Note that conditioned on the event $Z = z$, the random variables $X_{z}, Y$ are independent, and, furthermore, $Y$ is uniformly distributed over $\{0,1\}^d$. In addition, as $z \not\in B'$, $\minent{X_{z}} \ge k$. Thus, as $\ext$ is a strong $(k,\eps)$ seeded extractor,
\[
\left( \ext(X_{z},Y), Y \right) \approx_\eps \left( U_m, Y \right)\,.
\]
Let $Z'$ be the distribution obtained by sampling $z \sim Z$ conditioned on $z \notin B'$. Then,
\begin{align*}
T(Y,Y_1,\ldots,Y_t) &= \SD\big( \left(\ext(X,Y), Z , Y \right),  \left(U_m, Z, Y\right) \big) \\
&= \E_{z \sim Z} \left[ \SD\left( \left(\ext(X_{z},Y), Y \right),
          \left(U_m, Y \right)
   \right) \right] \\
&\le \Pr[Z \in B'] + \E_{z' \sim Z'}\left[ \SD\left( \left(\ext(X_{z'},Y), Y \right),
          \left(U_m, Y \right)
   \right) \right] \le 2\eps\,.
\end{align*}
By Markov's inequality, there exists a subset $B_1 \subseteq \{0,1\}^d$ of size $|B_1| \le 2\eps^{1/3} \cdot 2^d$ such that for every $y \notin B_1$,
\[
T(y,Y_1,\ldots,Y_t) \leq \eps^{2/3}\,.
\]
Thus, as $Y_1,\ldots, Y_t$ are independent, we get that for every $y \notin B_1$,
\[
\E_{(y_1,\ldots,y_t) \sim U_{td}}[T(y,y_1,\ldots,y_t)] \le \eps^{2/3}\,,
\]
concluding the proof of the claim.
\end{proof}

\begin{proposition}\label{c: pi construction}
There exist $(\widehat{A}_1,\ldots,\widehat{A}_t) \in  \advfampart^t_{\max(\Delta, \log t,1)} \cap \nc^t$ and $B_2 \subseteq \{0,1\}^d$ of size $\abs{B_2} \le 3\eps^{1/3}\cdot 2^d + t$ such that
\begin{enumerate}
  \item\label{it:c1} For every $y \in H$, $\widehat{A}_j(y) = \adv_j(y)$, and,
  \item\label{it:c2} For every $y \in \{0,1\}^d \setminus (H \cup B_2)$ it holds that $T\left(y,\widehat{A}_1(y), \ldots, \widehat{A}_t(y)\right) \le \eps^{1/3}$.
\end{enumerate}
\end{proposition}

\begin{proof}

Denote $\beta = 2^{-d} \cdot |B_1\setminus H|$. 
By \expref{Claim}{claim:average behavior}, $\beta \leq 2\eps^{1/3}$. Moreover, by Markov's inequality, for every $y \notin B_1$, there exists $B(y) \subseteq (\{0,1\}^d)^t$ of size at most $\abs{B(y)} \leq  \eps^{1/3} \cdot 2^{dt}$ such that for every $(y_1,\ldots,y_t) \notin B(y)$, \[T(y,y_1,\ldots,y_t) \le \eps^{1/3}\,.\]
Let $L = \{0,1\}^d \setminus (H \cup B_1)$. Fix an (arbitrary) ordering of the elements in $L$ and denote them by $y^1,y^2,\ldots,y^{\abs{L}}$. Let $c$ be the least integer larger than $\max(\eps^{1/3} \cdot 2^d, t)$, and set $\ell = \abs{L} - c$. Denote \[B_2 = (B_1\setminus H) \cup \{y^{\ell+1},\ldots, y^{\abs{L}}\}\,,\] observing that indeed
\[
|B_2| \le |B_1| + |L|-\ell = |B_1| +c \le 3\eps^{1/3}\cdot 2^d + t\,.
\]
%
We define a family of functions
\[
\set{\widehat{A}_{j,m} : 1 \le j \le t,\, 0 \le m \le \ell}\,,
\]
where $\widehat{A}_{j,0} \colon H  \rightarrow \{0,1\}^d$ and for $m \ge 1$,
\[\widehat{A}_{j,m} \colon H \cup \{y^1,\ldots, y^{m}\} \rightarrow \{0,1\}^d\,.\]
The above functions are constructed via the following algorithm which proceeds iteratively on $m$.%

%
\paragraph{The construction algorithm for the $\widehat{A}_{j,m}$ functions}
\begin{enumerate}
  	\item\label{it:assump1} For every $j \in [t]$, set $\widehat{A}_{j,0} = \adv_{j}|_{H}$. 
%
  	\item\label{it:assump2} For $m = 1,\ldots, \ell$, we will show in \expref{Claim}{c:step 2} below that there exists $(y^m_{1},\ldots,y^m_{t})\in (\{0,1\}^d)^t \setminus B(y^m)$ such that $y^m, y^m_{1},\ldots,y^m_{t}$ are pairwise distinct; furthermore, for every $j \in [t]$, $\abs{\widehat{A}_{j,m-1}^{-1}(y^m_{j})} < t$. Under these assumptions, for every $j \in [t]$ and $y \in H \cup \{y^1,\ldots, y^m\}$, set 
  	\[\widehat{A}_{j,m}(y) = \begin{cases}
  	y^m_{j} &  y = y^m; \\
  	\widehat{A}_{j,m-1}(y) & \text{otherwise}.
  	\end{cases}\]



\end{enumerate}

\begin{claim}\label{c:step 2}
	The underlying assumption of \expeqref{step}{it:assump2} in the algorithm above holds for every $m \in [\ell]$.
\end{claim}

\begin{proof}
We begin by setting  notation.
For every $m \in [\ell]$ and $j \in [t]$, we define the set $G_j^m \subseteq \{0,1\}^d \setminus \{ y^m\}$ of elements whose preimage is of size strictly less than $t$ with respect to the function $\widehat{A}_{j,m-1}$. Formally,
\[G_j^m = \set{y \in \{0,1\}^d \setminus \{ y^m\} : \abs{{\widehat{A}_{j,m-1}}^{-1}(y)} < t}\,.\]
With this notation, the hypothesis underlying \expeqref{step}{it:assump2} holds at iteration $m$ if there exists an element in $(G_1^m \times \cdots \times G_t^m) \setminus B(y^m)$ with pairwise distinct entries. To establish this, we start by bounding the size of $G_j^m$ from below for every fixed $j \in [t]$.

For an element $y$ not to be contained in $G^m_j$ there must be at least $t$ elements whose image under $\widehat{A}_{j,m-1}$ is $y$. At the beginning of the $m^{\text{th}}$ iteration, $m-1$ elements have been assigned an image at \expeqref{step}{it:assump2} and additional 
%
%
$|H| = \delta/2 \cdot 2^d$ 
%
elements were assigned at  \expeqref{Equation}{it:assump1}. Hence,
\begin{align}\label{eq:gm bound}
\abs{G_j^m} &\geq |\{0,1\}^d \setminus \{ y^m\}| - \frac{|H| + m-1}{t} \notag \\
&\ge 2^d -1 - \frac{|H| + \ell-1}{t}. %
%
\end{align}
As
$\ell = \abs{L} - c$ and since $L \cap H = \emptyset$ we have that
\[
|H|+\ell = |H| + \abs{L} - c \le |H \cup L| - c \le 2^d-c\,,
\]
and so
\[
\abs{G_j^m} \ge \left(1-\frac1{t}\right)2^d +  \frac{c+1}{t} -1.%
\]
Thus, as $y^m \not\in B_1$,
\begin{equation}\label{eq:pool}
\abs{\left(G_1^m \times \cdots \times G_t^m\right) \setminus B(y^m)} \ge \left(\left(1-\frac1{t} \right)2^d +  \frac{c+1}{t} -1\right)^t - \eps^{1/3} \cdot 2^d\,.
\end{equation}
%
%



Let $\ane \subseteq (\{0,1\}^d)^t$ be the largest set of vectors $v \in (\{0,1\}^d)^t$ such that $v_i \neq v_j$ for every pair of distinct $i,j \in [t]$. With this notation, to prove that the assumption underlying  \expeqref{step}{it:assump2} holds at step $m$, one must show that
\begin{equation}\label{eq:eqeqeqqeqeqseger}
\left(\left(G_1^m \times \cdots \times G_t^m \right)\setminus B(y^m)\right)\cap \ane \neq \emptyset\,.
\end{equation}
To this end, note that $\left|(\{0,1\}^d)^t \setminus \ane\right| \le \binom{t}{2} 2^{(t-1)d}$. Thus, by \expeqref{Equation}{eq:pool}, it suffices to show that
\begin{equation}\label{eq:long one}
\left(\left(1-\frac1{t} \right)2^d +  \frac{c+1}{t} -1\right)^t > \eps^{1/3} \cdot 2^d + \binom{t}{2} 2^{(t-1)d}\,. 
\end{equation}
For $t=1$, \expeqref{Equation}{eq:long one} is equivalent to
$
c > \eps^{1/3} \cdot 2^d
$
which readily follows by the definition of $c$. Consider then $t \ge 2$. By definition, $c \ge t$ and so %
%
$(c+1)/t-1 > 0$. 
%
Note further that $\left(1-\frac1{t}\right)^t \ge \frac1{4}$. Therefore, to satisfy \expeqref{Equation}{eq:long one}, it suffices to establish that
\[
\frac1{4} \cdot 2^{dt} \ge \eps^{1/3} \cdot 2^d + \binom{t}{2} 2^{(t-1)d}\,.
\]
It is straightforward to verify that the above equation follows per our assumption $d \ge 2\log{t} + 2$.
\end{proof}


We turn to extend the functions $\widehat{A}_{1,\ell}, \ldots, \widehat{A}_{t,\ell}$ to the domain $\{0,1\}^d$. To this end,
let $D = \{0,1\}^d \setminus (H \cup \{y^1,\ldots,y^\ell\})$ and denote its elements by $y^{\ell+1},\ldots,y^{\ell+e}$. For $m=1,\ldots,e$ we define the function \[
\widehat{\adv}_{j,\ell+m} \colon D \cup \{y^{\ell+1},\ldots, y^{\ell+m}\} \to \{0,1\}^d
\]
using the following iterative algorithm.

%
\paragraph{The algorithm for extending $\widehat{A}_{j,\ell}$ to $\widehat{A}_{j}$}
\begin{enumerate}
%
	%
	\item\label{It:assump21} For $m = 1,\ldots, e$, we will show in \expref{Claim}{c:step 3} below that there exist pairwise distinct $y^{\ell+m}_{1},\ldots,y^{\ell+m}_{t}\in \{0,1\}^d \setminus \{y^{\ell+m}\}$ such that for every $j \in [t]$, $\abs{\widehat{A}_{j,\ell+m-1}^{-1}(y^{\ell+m}_{j})} \le 1$. Under these assumptions, for every $j \in [t]$ and $y \in H \cup \{y^{1},\ldots, y^{\ell+m}\}$, set 
	\[\widehat{A}_{j,\ell+m}(y) = \begin{cases}
	y^{\ell+m}_{j} &  y = y^{\ell+m}; \\
	\widehat{A}_{j,\ell+m-1}(y) & \text{otherwise}.
	\end{cases}\]
	\item\label{It:assump22} For $j \in [t]$, set $\widehat{A}_{j} = \widehat{A}_{j,e}$.
	

	
\end{enumerate}

\begin{claim}\label{c:step 3}
	The underlying assumption in \expeqref{step}{It:assump21} of the algorithm above holds.
\end{claim}

\begin{proof}
	Fix $m \in [e]$. For every $j \in [t]$, the number of elements $z$ for which $|\widehat{A}^{-1}_{j,\ell+m-1}(z)| \ge 2$ is bounded above by $\frac1{2} \cdot 2^d$. Thus, when setting $\widehat{A}_{j,\ell+m}(y^{\ell+m})$ one has at least $2^d/2$ choices for an image with respect to this restriction. Recall that we also need to guarantee that
	\[
	y^{\ell+m}, \widehat{A}_{1,\ell+m}(y^{\ell+m}),\ldots,\widehat{A}_{t,\ell+m}(y^{\ell+m})\]
	are pairwise distinct. This can be achieved as $\frac1{2} \cdot 2^d+t+1 < 2^d$ per our assumption $d \ge 2\log{t}+2$.
\end{proof}

%
\paragraph{Analyzing the $\widehat{A}$ functions.}
First, note that the domain of each $\widehat{A}_1,\ldots,\widehat{A}_t$ is $\{0,1\}^d$. Recall that each of $A_1,\ldots,A_t$ has no fixed points by assumption and, furthermore, are non-colluding. Thus, $\widehat{A}_1,\ldots,\widehat{A}_t$ have no fixed points and are non-colluding when restricted to $H$. Moreover, by construction, the functions $\widehat{A}_1,\ldots,\widehat{A}_t$ are defined to have no fixed points and to be non-colluding outside of $H$ as well. Thus,  $(\mathrm{id},\widehat{A}_1,\ldots,\widehat{A}_t) \in \nc^{t+1}$, where $\mathrm{id} \colon \{0,1\}^d \to \{0,1\}^d$ is the identity function. 

We turn to show that $\widehat{A}_j \in \advfampart_{\max(\Delta, \log t,1)}$ for every $j \in [t]$.  
Recall that $\adv_j \in \advfampart_{\Delta}$. Both algorithms above assure that for any $j\in [t]$, and $y \notin \mathrm{Im}(\widehat{A}_j|_H)$ it holds that, $\abs{\widehat{A}_j^{-1}(y)} \leq \max(t,2).$
Observe that $f \in \advfampart_{\log t}$ if and only if 
for every $y \in \mathrm{Im}(f)$, $|f^{-1}(y)| \le t$,
and so
it holds that
\[
\widehat{A}_j \in \advfampart_{\Delta} \cup \advfampart_{\log \max(t,2)} =\advfampart_{\max(\Delta, \log t,1)}\,.
\]
This concludes the proof of \expref{Proposition}{c: pi construction}.
\end{proof}

We turn back to the proof of \expref{Theorem}{t: sp to nm}.
%
Define the random variables
\begin{align*}
    \mathcal{D}_{\widehat{A}} &= \mathcal{D}_{Y,\widehat{A}(Y)},\\
    \mathcal{I}_{\widehat{A}} &= \mathcal{I}_{\widehat{A}(Y)}\,,
\end{align*}
where $Y$ is uniformly distributed over $\{0,1\}^d$ and $\widehat{A}(Y) =  (\widehat{A}_1(Y),\ldots,\widehat{A}_t(Y))$.

%

\begin{claim}\label{claim:d is concentrated}
	There exists a set $B_3 \subseteq \{0,1\}^d$ of size $\abs{B_3} \le \sqrt{\eps} \cdot 2^d$ such that for every $y \in \{0,1\}^d \setminus B_3$,
	\[
	\SD\left(\mathcal{D}_{\widehat{A}}, \mathcal{D}_{y,\widehat{A}(y)}\right) \le \sqrt{\eps}.
	\]
\end{claim}

\begin{proof}
	%
	Denote 
	\begin{align*}
		\widehat{Z}(y) = \left(\ext(X,\widehat{A}_1(y)), \ldots, \ext(X,\widehat{A}_1(y))\right)\,.
	\end{align*}
	With this notation, 
	\begin{align*}
	\mathcal{D}_{\widehat{A}} &= \left( \ext(X,Y), \widehat{Z}(Y) \right)\,, \\
	\mathcal{D}_{y,\widehat{A}(y)} &= \left( \ext(X,y), \widehat{Z}(y) \right)\,.
	\end{align*}
%
	By \expref{Proposition}{c: pi construction}, $\widehat{A}_1,\ldots,\widehat{A}_t \in \advfampart_{\max(\Delta, \log t,1)}$; moreover, note that  $\mathrm{id} \in \advfampart_{\max(\Delta, \log t,1)}$. Observe that as $\ext$ is seed protecting against $\advfampart_{\max(\Delta, \log t,1)}^{t+1}$, and since $(\mathrm{id},\widehat{A}_1,\ldots,\widehat{A}_t) \in \nc^{t+1}$. it holds that %
	\begin{equation}\label{eq: low z}
	\E_{y\sim U_d}\big[\SD\big((\ext(X,y), \widehat{Z}(y)), (\ext(X,Y),\widehat{Z}(Y))\big)\big] \leq \eps\,.
	\end{equation}
	Let $B_3 \subseteq \{0,1\}^d$ be the set of all elements $y$ satisfying
	%
	\begin{align*}
	\SD\big((\ext(X,y), \widehat{Z}(y)),(\ext(X,Y), \widehat{Z}(Y))\big) \geq \sqrt{\eps}\,.
	\end{align*}
	By Markov's inequality, it follows that $\abs{B_3} \leq \sqrt{\eps}\cdot 2^d$, as stated.
%
	
\end{proof}

\begin{claim}\label{claim:i is concentrated}
	There exists a set
	 $B_4 \subseteq \{0,1\}^d$ of size $\abs{B_4} \le \sqrt{\eps} \cdot 2^d$ such that for every $y \in \{0,1\}^d \setminus B_4$,
	\[
	\SD\left(\mathcal{I}_{\widehat{A}}, \mathcal{I}_{\widehat{A}(y)}\right) \le \sqrt{\eps}\,.
	\]
\end{claim}
\begin{proof}
	Note that for every $y \in \{0,1\}^d$ it holds that  
	\[
	\SD\left(\mathcal{I}_{\widehat{A}}, \mathcal{I}_{\widehat{A}(y)}\right) = \SD\big(\widehat{Z}(Y)\,, \widehat{Z}(y)\big).
	\]
	By \expeqref{Equation}{eq: low z},
	\[
	\E_{y\sim U_d}\left[\SD\left(\mathcal{I}_{\widehat{A}}, \mathcal{I}_{\widehat{A}(y)}\right)\right] \leq \eps\,.  
	\]
	Let $B_4 \subseteq \{0,1\}^d$ be the set of all elements $y$ satisfying
	%
	\begin{align*}
	\SD\big(\mathcal{I}_{\widehat{A}}, \mathcal{I}_{\widehat{A}(y)}\big) \geq \sqrt{\eps}\,.
	\end{align*}
	By Markov's inequality, it follows that $\abs{B_4} \leq \sqrt{\eps}\cdot 2^d$, as stated.
\end{proof}


We are now ready to complete the proof. Write $B = B_1 \cup B_2 \cup B_3 \cup B_4$. Recall that \[\abs{H} = \frac{\delta}{2}\cdot 2^d = 7\eps^{1/3} \cdot 2^d\,.\]
By the above claims and using our hypothesis on $d$,
\[
\abs{B} \leq (5\eps^{1/3} + 2\sqrt{\eps})\cdot 2^d + t < 7\eps^{1/3} \cdot 2^d\,,
\]
and so there exists $y_h \in H \setminus B$.  On the other hand, 
\[
\abs{H \cup B} \le \left(\frac{\delta}{2} + 5\eps^{1/3} + 2\sqrt{\eps}  \right) 2^d + t < 2^d\,,
\]
where the last inequality follows as $\eps \leq 10^{-4}$ and, again, using our hypothesis on $d$. Hence, there exists $y_\ell \in \{0,1\}^d \setminus (H \cup B)$. 

By \expref{Claim}{claim:d is concentrated}, since $y_\ell,y_h \notin B_3$,
\begin{align*}
&\SD\left(\mathcal{D}_{y_h,\widehat{A}(y_h)}, \mathcal{D}_{\widehat{A}} \right) \le \sqrt{\eps}, \\
&\SD\left(\mathcal{D}_{y_\ell,\widehat{A}(y_\ell)}, \mathcal{D}_{\widehat{A}} \right) \le \sqrt{\eps}\,,
\end{align*}
and so
\begin{equation}\label{eq:dldh}
\SD\left(\mathcal{D}_{y_\ell,\widehat{A}(y_\ell)}, \mathcal{D}_{y_h,\widehat{A}(y_h)}\right) \le 2\sqrt{\eps}\,.
\end{equation}
By \expref{Claim}{claim:i is concentrated} and since $y_\ell,y_h \notin B_4$,
\begin{align*}
&\SD\left(\mathcal{I}_{\widehat{A}(y_h)}, \mathcal{I}_{\widehat{A}} \right) \le \sqrt{\eps}, \\
&\SD\left(\mathcal{I}_{\widehat{A}(y_\ell)}, \mathcal{I}_{\widehat{A}} \right) \le \sqrt{\eps}\,.
\end{align*}
Thus,
\begin{equation}\label{eq:ilih}
\SD\left(\mathcal{I}_{\widehat{A}(y_\ell)}, \mathcal{I}_{\widehat{A}(y_h)} \right) \le 2\sqrt{\eps}\,.
\end{equation}
Now, since $y_\ell \in \{0,1\}^d \setminus (H \cup B_2)$, by \expeqref{item}{it:c2} of \expref{Proposition}{c: pi construction},
\begin{equation}\label{eq:dlil}
\SD\left(\mathcal{D}_{y_\ell,\widehat{A}(y_\ell)}, \mathcal{I}_{\widehat{A}(y_\ell)}\right) \le \eps^{1/3}\,.
\end{equation}
By \expeqref{Equation}{eq:dldh}, \expeqref{Equation}{eq:ilih}, \expeqref{Equation}{eq:dlil}, and the triangle inequality,
\[
\SD\left(\mathcal{D}_{y_h,\widehat{A}(y_h)}, \mathcal{I}_{\widehat{A}(y_h)}\right) \le 5 \eps^{1/3}\,.
\]
However, $y_h \in H$ and so
\begin{align*}
\SD\left( \mathcal{D}_{y_h,\widehat{A}(y_h)}, \mathcal{I}_{\widehat{A}(y_h)} \right) &= T\left(y_h,\widehat{A}(y_h)\right) \\
&= T(y_h, \adv(y_h)) > \frac{\delta}{2}\,,
\end{align*}
contradicting our choice of $\delta$.
\end{proof}

Observe that for $\Delta = 0$ and $t = 1$, \expref{Theorem}{thm:warmup} shows that seed protection against
permutations is enough, and one does not need to devise an extractor against $\mathcal{F}_1$.
For a general $t > 1$ however, we suspect it is not true, and one needs to to handle $\mathcal{F}_\Delta$
for $\Delta > 0$ to get non-malleability against permutations. We note however, that if one is willing
to tolerate ``smoothness'' (in the sense of \expref{Definition}{def:smooth inv} and \expref{Definition}{def:smooth2} below), or a slight
error degradation, we can get non-malleability against permutations from seed-protecting extractors
against permutations, as long as $t$ is small enough. We omit the details.

\subsection{Colluding does not harm non-malleability}\label{sec:no nc}

\expref{Theorem}{t: sp to nm} established non-malleability against $\left(\advfampart_{\Delta} \cap \nofp\right)^t \cap \nc^t$.\footnote{We recall that $\mathcal{F}_{\Delta}$ is the set of functions $A \in \advfam_d$ for which $\me(A(U_d)) \ge d-\Delta$, $\nofp$ is the set of functions with no fixed points, and $\nc^t$ is
the set of $t$-tuples of non-colluding functions.}
However, in sharp contrast to seed-protecting extractors, colluding cannot help
adversaries in breaking non-malleable extractors. Intuitively this should be clear, as
redundant information should not help the adversary in distinguishing $\ext(X,Y)$
from uniform. Here we make it formal.

\begin{lemma}\label{lem:nm-with-coll} 
Let $\ext \colon \B^{n} \times \B^{d} \rightarrow \B^{m}$ be a $(k,\eps)$-non-malleable
extractor against $\mathcal{N}^{t} \cap \mathcal{X}^{t}$ (\ie, non-colluding functions
with no fixed points), so that $d \ge \log (t + 2)$. Then, $\ext$ is $(k,\eps)$-non-malleable
against $\mathcal{N}^{t}$ (\ie, a $t$-non-malleable extractor).
\end{lemma}
\begin{proof}
Let $A_1,\ldots,A_t \in \mathcal{A}_d \cap \mathcal{N}$ be any adversarial
functions with no fixed points. 
Fix an $(n,k)$-source $X$ and let $Y \sim \B^d$ be a uniform random variable,
independent of $X$.

We define the tuple of non-colluding functions
$A'_1,\ldots,A'_t$ as follows. Let $\mathbf{B} \subseteq \B^d$ be the set of $y$-s 
in which a colluding occurs. Namely, for each $y \in \mathbf{B}$ there exist distinct $i,j \in [t]$
for which $A_{i}(y)=A_{j}(y)$. Note that it is possible that different set of functions collude
separately, say $A_{i}(y)=A_{j}(y) = z$ and $A_{i'}(y)=A_{j'}(y) = z'$ for $z \neq z'$.
Given $y \in \B^d$, let $B(y) \subseteq [t]$ be the set of ``redundant'' 
adversaries for $y$. Formally, 
\[
B(y) = \set{i \in [t]: \text{there exists $j < i$ such that $A_{j}(y)=A_{i}(y)$}}\,.
\]
Note that if $y \notin \mathbf{B}$, $B(y)$ is empty. 
Also, given $y \in \mathbf{B}$, we denote by $I(y) = \set{A_{1}(y),\ldots,A_{t}(y)}$, and
take $E_1(y),\ldots,E_{t}(y)$ to be the first $t$ elements in $\B^{d} \setminus (I(y) \cup \{y\})$ in some fixed order.  
As $d \ge \log( t + 2)$, we can indeed do so. For every $i \in [t]$, we define
\[
A'_{i}(y) = 
\begin{cases}
A_{i}(y) & y \notin \mathbf{B}, \\
A_{i}(y) & y \in \mathbf{B} \wedge i \notin B(y), \\
E_{i}(y) & \text{otherwise}.
\end{cases}
\]
It then follows that $\left( A'_1,\ldots,A'_t \right) \in \mathcal{X}^t$. 

Denote $Z(y) = ( \ext(X,A_{1}(y)),\ldots,\ext(X,A_{t}(y)) )$, and likewise,
\[Z'(y) = ( \ext(X,A'_{1}(y)),\ldots,\ext(X,A'_{t}(y)) )\,.\] We further
define $Z_{\mathsf{reduce}}(y)$ to be joint distribution of the $\ext(X,A_i(y))$ for $i$-s which
are not in $B(y)$. Namely,
\[
Z_{\mathsf{reduce}}(y) = \bigcirc_{i \in [t] \setminus B(y)} \ext(X,A_{i}(y))
\]
where $\circ$ denotes concatenation. 
We record the following two easy claims.
\begin{claim}\label{claim:sec41}
For every $y \in \B^d$,
it holds that
\begin{align*}
&\SD((\ext(X,y),Z(y)),(U_m,Z(y))) =  \\ &\SD((\ext(X,y),Z_{\mathsf{reduce}}(y)),(U_m,Z_{\mathsf{reduce}}(y)))\,.
\end{align*}
\end{claim}
\begin{proof}
The claim follows from the following observation: For every three random variables $A$, $B$, and $C$,
it holds that $\SD((A,B,C,C),(U,B,C,C)) = \SD((A,B,C),(U,B,C))$, where $U$ is uniform over the support of
$A$ and independent of all other random variables.
\end{proof}

\begin{claim}\label{claim:sec42}
For every $y \in \B^d$ it holds that
\begin{align*}
&\SD((\ext(X,y),Z_{\mathsf{reduce}}(y)),(U_m,Z_{\mathsf{reduce}}(y))) \le \\
&\SD((\ext(X,y),Z'(y)),(U_m,Z'(y)))\,.
\end{align*}
\end{claim}
\begin{proof}
The claim readily follows from the data processing inequality, observing that
$Z'(y) = (Z_{\mathsf{reduce}(y)},A)$ for some random variable $A$.
\end{proof}

We can now finish the proof. As $\ext$ is non-malleable against non-colluding functions,
we know that
\[
\SD((\ext(X,Y),Z'(Y),Y),(U_m,Z'(Y),Y)) \le \eps\,.
\]
But
\[
\SD((\ext(X,Y),Z'(Y),Y),(U_m,Z'(Y),Y))  = \E_{y \sim U_d}\left[ \SD((\ext(X,y),Z'(y)),(U_m,Z'(y))) \right]\,,
\]
so
combining the above with \expref{Claim}{claim:sec41} and \expref{Claim}{claim:sec42}, it follows that
\begin{align*}
\SD((\ext(X,Y),Z(Y),Y),(U_m,Z(Y),Y)) &= \E_{y \sim U_d}\left[ \SD((\ext(X,y),Z(y)),(U_m,Z(y))) \right] \\
&= \E_{y \sim U_d}\left[ \SD((\ext(X,y),Z_{\mathsf{reduce}}(y)),(U_m,Z_{\mathsf{reduce}}(y))) \right] \\
&\le \E_{y \sim U_d}\left[ \SD((\ext(X,y),Z'(y)),(U_m,Z'(y))) \right] \\
&\le \eps\,,
\end{align*}
as desired. 
\end{proof}

%
%
%

The above lemma can be adapted to cases when the non-malleability is against
more restricted family of functions. In particular, we will need the following lemma.
\begin{lemma}\label{lem:nm-with-coll22} 
Let $\ext \colon \B^{n} \times \B^{d} \rightarrow \B^{m}$ be a $(k,\eps)$-non-malleable
extractor against $\left(\advfampart_{\Delta} \cap \nofp\right)^t \cap \nc^t$ for $\Delta  \ge 1$, so that $d \ge \log {t} + 3$. Then, $\ext$ is $(k,\eps)$-non-malleable
against $\left(\advfampart_{\Delta} \cap \nofp\right)^t$.
\end{lemma}
\begin{proof}
Inspecting the proof of \expref{Lemma}{lem:nm-with-coll}, we just need to make sure that
$A'_{1},\ldots,A'_{t}$ stay inside the family $\mathcal{F}_{\Delta}$. For every $i \in [t]$, let
$\mathbf{C}_i \subseteq \B^d$ be the set of seeds we re-wired. Namely,
\[
\mathbf{C}_i = \set{y \in \B^d : y \in \mathbf{B} \wedge i \in B(y)}\,.
\]
Recall that if $y \in \mathbf{C}_i$ then $A'_{i}(y) = E_{i}(y)$. We will modify the definition of 
$E_{i}$ to guarantee that $A'_i \in \mathcal{F}_{\Delta}$, while still satisfying
$(A'_1,\ldots,A'_t) \in \mathcal{X}^t \cap \nofp^t$.%

Let $\mathbf{G}_1(y) = \B^{d} \setminus (I(y) \cup \{y\})$ be the set of ``safe'' seeds. In \expref{Lemma}{lem:nm-with-coll}
we simply set $E_{i}(y)$ to be the $i$-th element of $\mathbf{G}_1(y)$. Now, we need to be just a bit
more careful. Let
\[
\mathbf{G}_2 = \set{z \in \B^d: \left| A_{i}^{-1}(z) \right| < 2^{\Delta}}\,.
\]
By an averaging argument, $|\mathbf{G}_2| \ge 2^{d-\Delta}$. Thus, 
\[
\left|\mathbf{G}_1(y) \cap \mathbf{G}_2\right| \ge 2^d - t - 1 - 2^{d-\Delta} \ge t\,,
\]
and we can set $E_{i}(y)$ to be the $i$-th element of $\mathbf{G}_1(y) \cap \mathbf{G}_2$.
Both properties now hold.
\end{proof}
%
%
%
%
%
%
%
%
%
%

Combining \expref{Therorem}{t: sp to nm} and \expref{Lemma}{lem:nm-with-coll22}, we get our main result.

\begin{corollary}\label{main:sec2}
Let $t \ge 1$ be an integer, and $\Delta \ge 1$.
Let $\ext\colon\B^{n}\times\B^{d}\rightarrow\B^m$ be a $(k,\eps)$-seeded extractor with $d \ge 2\log{t} + \log\frac1{\eps} + 2$ and $\eps \leq 10^{-4}$. Assume that $\ext$ is seed protecting against $\advfampart^{t+1}_{\max(\Delta, \log t)} \cap \nc^{t+1}$. Then, $\ext$ is non-malleable against $\left(\advfampart_{\Delta} \cap \nofp\right)^t$  for min-entropy $k+mt+\log\frac1{\eps}$ and error guarantee $14 \eps^{1/3}$.
\end{corollary}


%
\section{Seed-protecting   %
extractors from non-malleable extractors}

In this section, we show that non-malleability against permutations
implies seed protection for the class of permutations. We also prove a similar claim for a more
restricted kind of permutations -- $t$-cliques, which we define below. For the sake of generality, we consider the ``smooth'' variants of these classes.

\begin{definition}[smooth permutations]\label{def:smooth2}
Given $A \in \mathcal{A}_d$ and $\tau \in [0,1]$, we say that $A \in \Pi_\tau$ if there exists a set $G \subseteq \B^d$ 
with $|G| \ge (1-\tau) \cdot 2^d$ such that $A|_G$ is injective.
\end{definition}

To describe our next family of structured adversaries, we introduce the following notation. We say that a function $A \in \advfam$ is $\tau$-close to an involution if for all but $\tau$-fraction of $y \in \{0,1\}^d$ it holds that $A(A(y)) = y$ and $A(y) \neq y$. That is, the directed graph induced by the function $A$ has, but for one component of density $\tau$, an involution structure (\ie, a perfect matching). More generally, for $t \ge 1$, we formalize what it means for an adversarial function to be $\tau$-close to  $(t+1)$-cliques, or clusters. Fix functions $A_1,\ldots,A_t \in \advfam$. For $y\in \{0,1\}^d$, we define the neighborhood of $y$ by
	$
	\Gamma(y) = \left\{y,A_1(y),\ldots, A_t(y)\right\}.
	$	
	We say that $(A_1,\ldots,A_t) \in \MM_\tau^t$ if there exists $G \subseteq \{0,1\}^d$ with $\abs{G} \geq (1-\tau)\cdot2^d$ such that for every $y \in G$ it holds that
	$
	\Gamma(\Gamma(y)) = \Gamma(y),
	$
	and $|\Gamma(y)| = t+1$. So, intuitively, but for a density-$\tau$ component, the vertices are partitioned to cliques, or clusters, of size $t+1$. In fact, we consider a more structured variant, which we now formally define.


\begin{definition}[$t$-cliques]\label{def:cliques}
Given $A_1,\ldots,A_t \in \mathcal{A}_d$ and $p \ge t$, we say
$\left(A_1,\ldots,A_t \right) \in \mathcal{M}^{t}[p]$ if there exists a
partition $\B^d = C_1 \uplus \cdots \uplus C_\ell$, each $t \le |C_i| \le p$, such that
the following holds. For every $i \in [\ell]$ let $c_i = |C_i|$ and denote $C_i = \set{y^{0},\ldots,y^{c_i-1}}$, and $A_{1}\left(y^{j}\right) = z^{j}$. Then,
\begin{enumerate}
    \item $A_1$ restricted to $C_i$ is a permutation. That is, $\set{y^{0},\ldots,y^{c_i-1}} = \set{z^{0},\ldots,z^{c_i-1}}$. 
  %
    \item For any integers $j \in [c_i]$ and $2 \le r \le t$, $A_{r}\left( y^j \right) = z^{j+ r-1 \bmod c_i}$.
 %
\end{enumerate}
Note, in particular, that $\mathcal{M}^t[p] \subset \Pi^t$. For brevity, we denote $\MM^t[t] = \MM^t$.
\end{definition}

\begin{definition}[smooth $t$-cliques]\label{def:smooth inv}
	For positive integers $p \ge t$, and $\tau \in [0,1]$, we define the 
	set $\MM_\tau^t[p] \subseteq \advfam_d^t$ as follows. A tuple $
	\bm{A} =(A_1,\ldots,A_t) \in \MM_\tau^t[p]$ if there exists $G 
	\subseteq \{0,1\}^d$ with $\abs{G} \geq (1-\tau)\cdot2^d$ such that 
\[
\left( A_1|_G,\ldots,A_t|_G \right) \in \MM^t[p]
\]
where by the latter we mean formally that there exists a partition $G = C_1 \uplus \cdots \uplus C_\ell$, each $t \le |C_i| \le p$, such that
the conditions of \expref{Definition}{def:cliques} are met. For brevity, we denote $\MM_{\tau}^t[t] = \MM_\tau^t$.
%
\end{definition}

We begin with permutations, and recall that $\mathcal{X}^{t}$ denotes the set of $t$-tuples of non-colluding functions.

\begin{lemma}\label{lem:nm-to-sp}
Let $\ext \colon \B^{n} \times \B^{d} \rightarrow \B^{m}$ be a $(k,\eps)$-non-malleable extractor
against $\left( \Pi \cap \mathcal{N} \right)^t$ (\ie, permutations with no fixed points), and fix any $\tau \ge 0$. Then, $\ext$ is $(k,\eps')$-seed protecting against $\Pi^{t+1}_{\tau}  \cap \mathcal{X}^{t+1}$ for
$\eps' = 2(t+1)(\tau+\eps)$.
\end{lemma}
\begin{proof}
Let $A_1,\ldots,A_{t+1} \in \Pi_{\tau}$ be non-colluding. Let $X$ be an $(n,k)$-source, and let $Y \sim \B^{d}$ be 
a uniform random variable, independent of $X$. 
For each $i \in [t+1]$, let $G_i \subseteq \B^{d}$ be such that $A_{i}|_{G_i}$ is injective.
Writing $G = G_{1} \cap \cdots \cap G_{t+1}$, we have that $|G| \ge (1-(t+1)\tau) \cdot 2^{d}$.
%
For each $i \in [t+1]$, denote by $\widetilde{A}_{i} \in \Pi$ the permutation that is obtained by
keeping the function's value on $G$ and completing it to a permutation on $\B^d$ in such a way
that $\widetilde{A}_1,\ldots,\widetilde{A}_{t+1}$ do not collude. Observe that it is possible to do so as long as
$\left|\B^{d} \setminus G\right| \ge t+1$, which certainly holds.\footnote{Otherwise,
$\tau < 2^{-d}$ which implies $\tau = 0$ and we can simply take $\widetilde{A}_i = A_i$
for each $i \in [t+1]$.} 
%
In what follows, we denote
\begin{align*}
Z_1(Y) &= ( \ext(X,A_{1}(Y)),\ldots,\ext(X,A_{t+1}(Y)) ),\\
\widetilde{Z}_1(Y) &= ( \ext(X,\widetilde{A}_{1}(Y)),\ldots,\ext(X,\widetilde{A}_{t+1}(Y)) )\,.
\end{align*}
First, note that
\begin{equation}\label{eq:sec41}
\SD((Y,Z_1(Y)),(Y,\widetilde{Z}_1(Y))) \le \Pr[Y \notin G] \le (t+1)\tau\,.
\end{equation}
Next, as $\widetilde{A}_{1}^{-1}$ is a permutation, $\widetilde{A}_{1}^{-1}(Y)$  distributes the same as $Y$ does. Thus,
\begin{equation}\label{eq:sec42}
\SD((U_d,\widetilde{Z}_1(Y)),(Y,\widetilde{Z}_1(Y))) = 
\SD((U_d,\widetilde{Z}_1(\widetilde{A}_{1}^{-1}(Y))),(\widetilde{A}_{1}^{-1}(Y),\widetilde{Z}_1(\widetilde{A}_{1}^{-1}(Y))))\,.
\end{equation}
By the data-processing inequality,
we can apply $\widetilde{A}_1$ on the prefix of both random variables without increasing
the statistical distance; \expeqref{Equation}{eq:sec42} becomes
\begin{equation}\label{eq:sec423}
\SD((U_d,\widetilde{Z}_1(Y)),(Y,\widetilde{Z}_1(Y))) = 
\SD((U_d,\widetilde{Z}_1(\widetilde{A}_{1}^{-1}(Y))),(Y,\widetilde{Z}_1(\widetilde{A}_{1}^{-1}(Y))))\,.
\end{equation}
For $i \in [t]$, define $P_i = \widetilde{A}_{i+1} \circ \widetilde{A}_{1}^{-1}$. Then,
\begin{align*}
\widetilde{Z}_1(\widetilde{A}_{1}^{-1}(Y)) &= \left( \ext(X,Y), \ext\left(X,\widetilde{A}_{2}(\widetilde{A}_{1}^{-1}(Y))\right), \ldots,\ext\left(X,\widetilde{A}_{t+1}(\widetilde{A}_{1}^{-1}(Y))\right)\right) \\ &=
\left( \ext\left( X,Y \right),\ext\left( X,P_{1}(Y) \right),\ldots,\ext\left( X,P_{t}(Y) \right) \right)\,.
\end{align*}
Note that for every $i \in [t]$, $P_i$ is a permutation. Furthermore, observe that as $(\widetilde{A}_1,\ldots,\widetilde{A}_{t+1}) \in \nc^{t+1}$,  $P_i$ has no fixed points. Indeed, assuming
towards a contradiction that $P_i$ satisfies $P_i(y)=y$ for some $y \in \B^{d}$, we get
that $\widetilde{A}_{1}^{-1}(y) = \widetilde{A}^{-1}_{i+1}(y)$, which is impossible since
$\widetilde{A}_{1}(z) \neq \widetilde{A}_{i+1}(z)$ for any $z \in \B^{d}$.
By the non-malleability of $\ext$, we know that
\[
\SD((Y,\widetilde{Z}_1(\widetilde{A}_1^{-1}(Y))),(Y,U_m,Z_{2}(Y))) \le \eps\,,
\]
for $Z_{2}(Y) = (\ext(X,P_{1}(Y)),\ldots,\ext(X,P_{t}(Y)))$, and the same is true without conditioning on $Y$.
Thus, by the triangle inequality, we get
\begin{align}\label{eq:sec44}
\SD((U_d,\widetilde{Z}_1(\widetilde{A}_{1}^{-1}(Y))),(Y,\widetilde{Z}_1(\widetilde{A}_{1}^{-1}(Y)))) &\le 
\SD((U_d,\widetilde{Z}_1(\widetilde{A}_{1}^{-1}(Y))),(U_d,U_m,Z_{2}(Y))) + \nonumber \\
& ~~~~\SD((U_d,U_m,Z_{2}(Y)),(Y,U_m,Z_{2}(Y))) + \nonumber \\
& ~~~~\SD((Y,U_m,Z_{2}(Y)),(Y,\widetilde{Z}_1(\widetilde{A}_{1}^{-1}(Y)))) \nonumber \\
&\le 2\eps + \SD((U_d,U_m,Z_{2}(Y)),(Y,U_m,Z_{2}(Y)))\,.
\end{align}

We continue in the same manner. Observing that $P_{1}^{-1}(Y)$ distributes the same as $Y$
does, and using the data-processing inequality, we have
\begin{equation}\label{eq:sec45}
\SD((U_{d+m},Z_{2}(Y)),(Y,U_m,Z_{2}(Y))) = 
\SD((U_{d+m},Z_{2}(P_{1}^{-1}(Y))),(Y,U_m,Z_{2}(P_{1}^{-1}(Y))))\,.
\end{equation}
For $i \in [t-1]$ define the permutation $Q_i = P_{i+1} \circ P_1^{-1}$.
Similarly to the above argument,
\begin{align*}
Z_{2}(P_{1}^{-1}(Y)) &= \left( \ext(X,Y),\ext(X,P_{2}(P_{1}^{-1}(Y))),\ldots,\ext(X,P_{t}(P_{1}^{-1}(Y))) \right) \\
&= \left( \ext(X,Y),\ext(X,Q_{1}(Y)),\ldots,\ext(X,Q_{t-1}(Y)) \right)\,.
\end{align*}
This time, the fact that every $Q_i$ has no
fixed points follows from the fact that $\widetilde{A}_{2}(z) \neq \widetilde{A}_{i+2}(z)$ for
any $z \in \B^{d}$. Using the non-malleability, together with
\expeqref{Equations}{eq:sec44}, \expeqref{Equations}{eq:sec423},
and \expeqref{Equations}{eq:sec45}, we get 
\begin{align*}
\SD((U_d,\widetilde{Z}_1(Y)),(Y,\widetilde{Z}_1(Y))) &\le 2\eps + \SD((U_{d},U_{m},Z_{2}(Y)),(U_{d},U_{2m},Z_{3}(Y))) + \\
&~~~~ \SD((U_d,U_{2m},Z_{3}(Y)),(Y,U_{2m},Z_{3}(Y))) + \\
&~~~~ \SD((Y,U_{2m},Z_{3}(Y)),(Y,U_{m},Z_{2}(Y))) \\
&\le 4\eps + \SD((U_d,U_{2m},Z_{3}(Y)),(Y,U_{2m},Z_{3}(Y)))\,,
\end{align*}
for $Z_{3}(Y) = \left( \ext(X,Q_{1}(Y)),\ldots,\ext(X,Q_{t-1}(Y)) \right)$. We continue
this process inductively, and eventually obtain
\begin{equation}\label{eq:sec46}
\SD((U_d,\widetilde{Z}_1(Y)),(Y,\widetilde{Z}_1(Y))) \le 2(t+1)\eps + \SD((U_d,U_{(t+1)m}),(Y,U_{(t+1)m})) = 2(t+1)\eps\,.
\end{equation}
Combining \expeqref{Equations}{eq:sec46} and 
\expeqref{Equation}{eq:sec41}, we get
\begin{align*}
\SD((U_d,Z_1(Y)),(Y,Z_1(Y))) &\le \SD((Y,Z_1(Y)),(Y,\widetilde{Z}_{1}(Y))) + \\ &~~~~\SD((Y,\widetilde{Z}_1(Y)),(U_d,\widetilde{Z}_1(Y))) + \\
&~~~~\SD((U_d,\widetilde{Z}_1(Y)),(U_d,Z_{1}(Y))) \\
&\le (t+1)(2\eps+\tau) + \SD(\widetilde{Z}_1(Y),Z_{1}(Y))\,.
\end{align*}
To bound the last term of the above inequality, note that
\[
\SD(\widetilde{Z}_1(Y),Z_{1}(Y)) \le \SD((Y,\widetilde{Z}_1(Y)),(Y,Z_{1}(Y))) \le (t+1)\tau\,,
\]
where the last inequality follows by \expeqref{Equation}{eq:sec41}. This concludes the proof.
\end{proof}

Next, we prove a similar lemma for $t$-cliques. 
\begin{lemma}\label{lem:nm-to-sp2}
Let $\ext \colon \B^{n} \times \B^{d} \rightarrow \B^{m}$ be a $(k,\eps)$-non-malleable extractor
against $\MM^t \cap \mathcal{N}^{t}$, and fix any $\tau \ge 0$. Then, $\ext$ is a $(k,\eps')$-seed-protecting extractor against $\MM^{t+1}_{\tau} \cap \nc^{t+1}$ for
$\eps' = 2(t+1)\eps + 2\tau$. 
\end{lemma}
\begin{proof}
Given %
non-colluding  %
$\bm{A} = \left(A_1,\ldots,A_{t+1} \right) \in \MM_{\tau}^{t+1}$,
an $(n,k)$-source $X$ and a uniform $Y \sim \B^d$ independent of $X$,
the proof proceeds similarly to \expref{Lemma}{lem:nm-to-sp}. Let $G \subseteq \{0,1\}^d$ be the set defined in \expref{Definition}{def:smooth inv} with respect to $(A_1,\ldots,A_{t+1}) \in \MM^{t+1}_{\tau}$, and recall that $|G| \le \tau \cdot 2^d$.
We define 
%
%
\[\widetilde{\bm{A}} = \left(\widetilde{A}_1,\ldots,\widetilde{A}_{t+1}\right) \in \MM^{t+1}\]
by setting $\widetilde{A}_i$ to agree with $A_i$ on $G$ for every $i$.
%
Completing $\widetilde{A}_i|_{G}$ to functions on $\B^d$ while maintaining the $\MM^{t+1}$ property can
be done by arbitrarily selecting $t+1$ inputs that were not assigned yet, iteratively, and assigning them
to form a clique.\footnote{Formally, choose $y_1,\ldots,y_{t+1}$ that
were not assigned by $\widetilde{A}_1$ yet, namely, with no preimage in $\widetilde{A}_1|_G$. 
Assign $y_1 \rightarrow y_1,\ldots,y_{t+1} \rightarrow y_{t+1}$ in $\widetilde{A}_1$, 
$y_1 \rightarrow y_2,y_2 \rightarrow y_3,\ldots,y_{t+1}\rightarrow y_{1}$ in $\widetilde{A}_2$,
and so on. Note that when $t+1$ does not divide $2^d$, we can make some cliques larger. We do not address this issue formally and it does not affect the statement.}
Note that by the definition of $\MM^{t+1}$, the functions in $\widetilde{\bm{A}}$ do not collude.  


Define, inductively, the following set of functions. 
For $i \in [t+1]$, $P_i^{(0)} = \widetilde{A}_i$. For every
$j \in [t]$ and $i \in [t+1-j]$, we define
\begin{equation}\label{eq:composing perm}
P_i^{(j)} = P_{i+1}^{(j-1)} \circ \left( P_{1}^{(j-1)} \right)^{-1}\,.
\end{equation}
Next, define, exactly as in \expref{Lemma}{lem:nm-to-sp},
\[
\widetilde{Z}_1(Y) = \left( \ext\left(X,P_{1}^{(0)}(Y)\right),\ldots, \ext\left(X,P_{t+1}^{(0)}(Y)\right) \right)\,.
\]
Moreover, for every $j \in [t+1]$, we define
\[
Z_{j+1}(Y) = \left( \ext\left(X,P_{1}^{(j)}(Y)\right),\ldots, \ext\left(X,P_{t-j+1}^{(j)}(Y)\right)\right)\,.
\]
The crux of the proof is establishing the following inequality for every $j \ge 2$.
\begin{align}\label{eq:pm}
\SD((U_d,\widetilde{Z}_1(Y)),(Y,\widetilde{Z}_1(Y))) &\le 2(j-1)\eps + \nonumber \\
& ~~~~\SD((U_{d},U_{(j-1)m},Z_j(Y)),(Y,U_{(j-1)m},Z_{j}(Y)))\,.
\end{align}
Following the same reasoning as in \expref{Lemma}{lem:nm-to-sp}, \expeqref{Equation}{eq:pm} holds if the 
following conditions are met:
\begin{enumerate}
\item\label{it:pm1}	For every $j \ge 0$, $\left(P_{1}^{(j)}\right)^{-1}(Y)$ distributes the same as $Y$ does.
\item\label{it:pm2}	For every $j \ge 1$ and $i \le t+1-j$, $P_{i}^{(j)}$ has no fixed points.
\item\label{it:pm3}	For every $j \ge 0$ it holds that $( P_{1}^{(j)},\ldots,P_{t-j+1}^{(j)} ) \in \MM^{t-j+1}[t+1]$.
Note that a non-malleable extractor against $\MM^{t+1}$
is also non-malleable against $\MM^{t'}[t+1]$ for any $t' \le t+1$. 
\end{enumerate}

\expeqref{Item}{it:pm1} readily holds, since each $\widetilde{A}_i$ is a permutation, and permutations are closed
under inversion and composition. To see that \expeqref{item}{it:pm2} holds, fix some
$j \in [t]$ and $i \in [t+1-j]$ and consider $P_{i}^{(j)}$. Assume
towards a contradiction that $P_{i}^{(j)}(y) = y$ for some $y \in \B^d$. Thus,
\[
P_{i+1}^{(j-1)} \left( \left(P_{1}^{(j-1)}\right)^{-1}(y)\right) = y\,,
\]
so
\[
\left(P_{1}^{(j-1)}\right)^{-1}(y) = \left( P_{i+1}^{(j-1)} \right)^{-1}(y)\,,
\]
which means there exists $z \in \B^d$ such that $P_{i+1}^{(j-1)}(z) = P_{1}^{(j-1)}(z)$.
But due to \expeqref{item}{it:pm3} which we now prove, $\left( P_{1}^{(j-1)},\ldots,P_{t-j+2}^{(j-1)} \right) \in \MM^{t-j+2}$,
so in particular they do not collude, so we have a contradiction. Indeed, what is left is 
to prove \expeqref{item}{it:pm3}.
\begin{claim}
For every $j \ge 0$ it holds that $( P_{1}^{(j)},\ldots,P_{t-j+1}^{(j)} ) \in \MM^{t-j+1}[t+1]$.
\end{claim}
\begin{proof}
Starting from $(t+1)$-cliques in the $0$-th level, in general, the $j$-th level
will also form $(t+1)$-cliques. We will prove this by induction on $j$. For $j = 0$, it follows
by our construction. Assuming that it holds for some $j \ge 0$, we inspect the $(j+1)$-th level.

We characterize the cliques in the $j$-th level as follows. We can partition the domain (and codomain)
to $\B^d = C_1 \uplus \cdots \uplus C_{\ell}$. By definition, a full characterization of
$P_1^{(j)},\ldots,P_{t-j+1}^{(j)}$ can be given by a permutation
$\phi_i \colon C_i \rightarrow C_i$, for each $i \in [\ell]$. Indeed, 
for any $y \in C_i$ for some $i \in [\ell]$, $P_{1}^{(j)}(y) = \phi_{i}(y)$, and
for $r \ge 2$, $P_{r}^{(j)} = \phi_{i}^r$. Namely, $\phi_i$ generates the permutations $P_1^{(j)},\ldots,P_{t-j+1}^{(j)}$ restricted to $C_i$.
We claim that this structure is preserved for the next level, with the same clique structure. Indeed, for each $i \in [\ell]$ observe that the permutation ${\phi}_i$ is also a generator for $P_1^{(j+1)},\ldots,P_{t-j}^{(j+1)}$ restricted to $C_i$ as, by \expeqref{Equation}{eq:composing perm} applied with $i=1$,
\[
P^{(j+1)}_{1} = P_{2}^{(j)} \circ \left( P_{1}^{(j)} \right)^{-1} = \phi_i^2 \circ \phi_i^{-1} = \phi_i\,.
\]
%
%
%
Recalling that $P^{(j+1)}_{r} = P_{r+1}^{(j)} \circ \left( P_{1}^{(j)} \right)^{-1}$, we see that indeed
computing $P^{(j+1)}_{r}(y)$ amounts to finding the $i$ for which $y \in C_i$
and computing $\phi_i^r(y)$. Thus, the $(j+1)$-th level belong to $\MM^{t-j}[t+1]$.
\end{proof}
\end{proof}


\section{1-seed protecting and two-source extractors}\label{sec:two source}

In this section we prove \expref{Theorem}{thm:two source sp intro}. We prove each direction separately in \expref{Lemma}{lem: sp is ts} and \expref{Lemma}{lem:two-source} below. Recall that $\mathcal{F}_{\Delta}$ is the set of functions $A \in \advfam_d$ for which $\me(A(U_d)) \ge d-\Delta$.

\begin{lemma}\label{lem: sp is ts}
Let  $\ext \colon \B^{n} \times \B^{d} \rightarrow \B^{m}$ be a $(k_1,\eps)$ $1$-seed-protecting extractor against $\advfampart_{d-k_2}$ for $k_2\leq d-1$. Then, $\ext$ is a $(k_1,k_2, 3\eps)$ two-source extractor.
\end{lemma}
%

\begin{proof}
 Assume towards a contradiction that $\ext$ is not such a two-source extractor, and let $X_1$,$ X_2$ with  $\minent{X_1} \geq k_1$ and $\minent{X_2} \geq k_2$ be such that 
	\[
	\SD\big(\ext(X_1,X_2), U_m \big) > 3\eps\,,
	\]
and we assume without loss of generality that both $X_1$ and $X_2$
are flat.\footnote{A $k$-source is \emph{flat} if it is uniformly distributed over a set of size $2^k$. It is well-known that one can assume $X_1$ and $X_2$ are flat, since any $k$-source is a convex combination of flat $k$-sources~\cite{CG88}.}
	Let $Y$ be the uniform distribution over $\{0,1\}^d$, independent of $(X_1, X_2)$.
	Recall that $\ext$ is a $(k_1, \eps)$ extractor, so 
	\[
	\SD(\ext(X_1,Y) ,U_m) \le \eps.
	\]
Order the $y$-s according to $\SD(\ext(X_1,y),U_m)$, and let $G_1 \subseteq \B^d$
be the $2^{d-1}$ bottom ones (\ie, for which $\SD(\ext(X_1,y),U_m)$ is smaller).
As $G_1$ can be indexed using $d-1$ bits, there exists an injection
$\adv_0 \colon \B^{d-1} \rightarrow \B^{d}$ that maps to $G_1$ uniformly, and satisfies
		\begin{equation}\label{eq:1}
	\SD(\ext(X_1,\adv_0(Y')),U_m) \leq \eps\,,
	\end{equation}
	where $Y'$ is the uniform distribution over $\{0,1\}^{d-1}$, independent of
	all other variables.

Denote $G_2 = \mathrm{Supp}(X_2)$ and recall that $|G_2| = 2^{k_2}$.
Again, since $k_2 \leq d -1$, we can define an injection $\adv_1 \colon \{0,1\}^{d-1} \rightarrow \{0,1\}^d$
so that $\adv_1(U_{d-1}) = U_{G_2}$, and then
	\begin{equation}\label{eq:2}
	\SD(\ext(X_1,\adv_1(Y')),U_m) > 3\eps\,.
	\end{equation}
Next, define $\adv \in \advfam_d$ such that \[\adv(y) = \adv_{y_1}(y_{[2:n]})\,.\]
\begin{claim}
It holds that $\adv\in \advfampart_{d-k_2}$.
\end{claim}
\begin{proof}
Let $Z = A(U_d)$. Clearly, $\mathrm{Supp}(Z) \subseteq G_1 \cup G_2$.
If $z \in G_1$,
\begin{align*}
\Pr_{y \sim U_d}[A(y) = z] &= \frac{1}{2} \cdot \Pr_{y \sim U_d}[\adv_0(y_{[2,n]})  = z]
+ \frac{1}{2} \cdot \Pr_{y \sim U_d}[\adv_1(y_{[2,n]})  = z] \\
&= \frac{1}{2} \cdot \frac{1}{|G_1|} + \frac{1}{2} \cdot \frac{\mathbf{1}_{z \in G_2}}{|G_2|}
\le 2^{-k_2}\,.
\end{align*}
The same bound applies for $z \in G_2$ in a similar manner.
\end{proof}


We will now show that using $\adv$, an adversary can learn the first bit of the seed, in contradiction to the fact that $\ext$ is seed protecting.
Define \[R_{Y_1} \triangleq \ext(X_1,\adv_{Y_1}(\lseed))\,,\] and note that by \expeqref{Equation}{eq:1}, \expeqref{Equation}{eq:2} and the triangle inequality, it holds that 
$\SD(R_0,R_1) > 2\eps.$
Then,
\begin{align*}
\SD(( Y,\ext(X_1,\adv(Y))), (U_d,\ext(X_1,\adv(Y)))) &\geq
\SD(( Y_1,\ext(X_1,\adv(Y))), (U_1, \ext(X_1,\adv(Y)))) \\ 
&=\SD(( Y_1, R_{Y_1}), (U_1,R_{Y_1}))\,.
\end{align*}
%
%
%
Finally, observe that 
\begin{align*}
\SD(( Y_1, R_{Y_1}), (U_1,R_{Y_1}) &= \frac{1}{2}\sum_{r}\left| \Pr[Y_1 = 1 \wedge R_{1} = r] - \frac{1}{2}\Pr[R_{Y_1}=r] \right| \\
&~~~+ \frac{1}{2}\sum_{r}\left| \Pr[Y_1 = 0 \wedge R_{0} = r] - \frac{1}{2}\Pr[R_{Y_1}=r] \right| \\ 
&= \frac{1}{4}\sum_{r}\left| \Pr[R_0 = r] - \Pr[R_1=r] \right| \\ &= \frac{1}{2} \cdot \SD(R_0,R_1) \\ &> \eps\,,
\end{align*}
which is a contradiction to that $\ext$ is 1-seed protecting per our hypothesis.%



\end{proof}

The other direction also holds.

\begin{lemma}\label{lem:two-source}
Let $\ext \colon \B^n \times \B^d \rightarrow \B^m$ be a 
$(k_1,k_2,\eps)$ two-source extractor which is strong in the second source. Then, $\ext$ is a $(k_1,2\eps)$
$1$-seed protecting against $\advfampart_{d-k_2}$.
\end{lemma}

\begin{proof}
Assume towards a contradiction that $\ext$ is not $1$-seed protecting,
so there exists an $(n,k_1)$-source $X_1$ and $\adv \in \advfampart_{d-k_2}$
such that
	\[
	\SD((\ext(X_1,\adv(Y)), Y), (\ext(X_1,\adv(Y)), Y')) > 2\eps\,,
	\]
where $Y, Y' \sim \B^d$ are uniform and independent random variables, also
independent of $X_1$.
	By the triangle inequality, it follows that at least one of the following holds:
\begin{enumerate}
\item\label{it:tw-source1} $\SD((\ext((X_1,A(Y)), Y'), (U_{m}, Y)) > \eps$,
\item $\SD((\ext(X_1,A(Y)), Y), (U_{m}, Y)) > \eps$.
\end{enumerate}	
First assume the first inequality holds. As $Y'$ is independent of $X_1$ and $Y$,
\begin{align*}
\SD((\ext((X_1,A(Y)), Y'), (U_{m}, Y)) &= 
\SD(\ext(X_1,A(Y)) \times Y', U_m \times Y) \\
&= \SD(\ext(X_1,A(Y)) \times Y', U_m \times Y') \\ &= \SD(\ext(X_1,A(Y)) , U_m)\,.
\end{align*}
As $\adv \in \advfampart_{d-k_2}$, $\minent{\adv(Y)} \ge k_2$ and by the fact that $\ext$ is a two-source extractor, it follows that  $\SD(\ext(X_1,A(Y)) , U_m) \le \eps$,
contradicting \expeqref{item}{it:tw-source1}.
	
Next, assume that the second inequality holds.
We have that 
\begin{align*}
	\SD((\ext(X_1,\adv(Y)), \adv(Y)), (U_{m}, \adv(Y))) &= \E_{z \sim \adv(Y)}\left[ \SD(\ext(X_1,z),U_{m}) \right] \\
	&= \E_{y \sim Y}\left[ \SD(\ext(X_1,\adv(y)),U_{m}) \right] \\
	&=\SD((\ext(X_1,\adv(Y)), Y), (U_{m}, Y)) > \eps\,.
\end{align*}
But since $\ext$ is a strong $(k_1,k_2,\eps)$ two-source extractor,
\[\SD((\ext(X_1,\adv(Y)), \adv(Y)), (U_{m}, \adv(Y))) \le \eps\,,\] in contradiction.
\end{proof}

\subsection{Lower bounds for 1-seed-protecting extractors}

%
In light of the above result, lower bounds for (unbalanced) two-source 
extractors imply lower bounds for $1$-seed-protecting extractors.
We use the following standard lower bound.\footnote{The lower bound for
$m = 1$ and any nontrivial $\eps$ follows from bounds on strong dispersers
and their connection to Ramsey graphs \cite{RTS00,BDT20}. The entropy loss
and the $2\log(1/\eps)$ factor in $d$ and $k$ follows from lower bounds
on strong seeded extractors~\cite{RTS00}.}
\begin{theorem}
Let $\ext \colon \B^{n} \times \B^{d} \rightarrow \B^{m}$ be a 
$(k_1,k_2,\eps)$ two-source extractor which is strong in the second source. Then,
$m \le k_1 - 2\log(1/\eps)+O(1)$, $k_2 \ge \log(n-k_1)+2\log(1/\eps)-O(1)$
and $k_1 \ge \log(d-k_2)+2\log(1/\eps)-O(1)$.
\end{theorem}

We can thus conclude:
\begin{corollary}
Let $\ext \colon \B^{n} \times \B^{d} \rightarrow \B^{m}$ be a $(k,\eps)$ $1$-seed protecting against $\advfampart_{\Delta}$ for some $\Delta > 0$.
Then, 
$m \le k - 2\log(1/\eps)+O(1)$, $d \ge  \log(n-k)+\Delta+2\log(1/\eps)-O(1)$
and $k \ge \log\Delta+2\log(1/\eps)-O(1)$.
\end{corollary}

The fact that no $1$-seed-protecting extractors against $\advfampart_\Delta$
exist for $\Delta$ approaching $d$ can be established in a more straightforward
way. 
Indeed, in \expref{Claim}{c:lower bound image two intro} we showed why
$\Delta = d-1$ is unattainable.


%
\section{Non-explicit 
\texorpdfstring{$t$-seed-protecting}{t-seed-protecting} extractors}\label{sec:nexp}

In this section we prove the existence of a seed-protecting extractor against non-colluding, entropy-preserving, adversaries via a probabilistic argument.

\begin{theorem}\label{thm:non-explicit}
Let $n,k,m,d,t \in \mathbb{N}$, $\Delta \ge 0$ and $\eps > 0$ be such that
\[
k \ge tm+2\Delta+2\log\frac{1}{\eps}+O(\log d + \log t)\,.
\]
Then, there exists a $(k,\eps)$-seed-protecting extractor $\ext \colon \B^{n} \times \B^{d} \rightarrow \B^{m}$ against
$\mathcal{F}_{\Delta}^{t} \cap \nc^t$ with
\[
d  = \log(n-k) + 2\Delta +2\log\frac{1}{\eps}+O(\log d + \log t)\,.
\]
\end{theorem}

%
\begin{remark}  %
\expref{Theorem}{thm:non-explicit} tells us we cannot take $\Delta$ to be larger than %
%
$d/2$, 
%
meaning
we are at the \emph{weak seeds} regime. In contrast, for $t = 1$, we know from the
equivalence to two-source extractors (see \expref{Lemma}{lem:two-source}) that we can take
$\Delta$ to be much larger, roughly $d-\log n$. An interesting open problem is 
whether there is a real barrier going from $t = 1$ to larger $t$-s or whether it is a mere
artifact of our proof.
\end{remark}   %

\begin{proof}
Choose $\ext \colon \B^{n} \times \B^{d} \rightarrow \B^{m}$ uniformly at random.
Fix an $(n,k)$ source $X$, and we assume without loss of generality that it is flat. Also, fix non-colluding functions $\adv_1,\ldots,\adv_t \in \advfampart_{\Delta}$. Let $Y \sim \B^{d}$ be uniform and independent of $X$.
Write $N = 2^n$, $K = 2^k$, $D = 2^d$, and
\[
Z(x,y) = \left( \ext(x,\adv_1(y)),\ldots,\ext(x,\adv_t(y)) \right)\,,
\]
observing that for every fixed $x$ and $y$, $Z$ is a random variable whose
randomness comes from $\ext$. We want to bound the probability that
\[
\SD( (Y,Z(X,Y)),(U_d,Z(X,Y)) ) > \eps\,.
\]
Fix $T \colon \B^{d+mt}\rightarrow\B$ and for each $y \in [D]$, denote by
$T_y \colon \B^{mt} \rightarrow \B$ the corresponding restriction of $T$.
Then, we want to bound the probability over $\ext$ that
\[
\E_{w \sim [D]}\left[ \E_{X}\left[ T_{w}(Z(X,w)) \right] \right] - \E_{w \sim [D]}\left[ \E_{X,Y}\left[ T_{w}(Z(X,Y)) \right] \right] > \eps\,.
\]
Write the expression on the left hand side as
\[
\E_{x \sim X,w \sim [D]}\left[ T_{w}(Z(x,w)) - \E_{Y}\left[ T_{w}(Z(x,Y)) \right] \right]\,,
\]
and define
\[
Q(x,w) = T_{w}(Z(x,w)) - \E_{Y}\left[ T_{w}(Z(x,Y)) \right] \,.
\]
First, we argue,
\begin{claim}
For any $x \in \B^n$ and $w \in [D]$ it holds that $\E[Q(x,w)] = 0$.
\end{claim}
\begin{proof}
By our assumption on $\adv_1,\ldots,\adv_t$, the values $\adv_{1}(w),\ldots,\adv_{t}(w)$ 
are distinct, so $Z(x,w)$ is uniform over $\B^{mt}$, and thus $\E[T_{w}(Z(x,w))] = \mu(T_w)$.
The claim now follows from the linearity of expectation.
\end{proof}

Write
\[
Q(x,w) =\left( T_{w}(Z(x,w)) - \mu(T_w) \right) - \left( \E_{Y}\left[ T_{w}(Z(x,Y)) \right] - \mu(T_w) \right) 
\triangleq Q_{1}(x,w) - Q_{2}(x,w)\,,
\]
each $Q_{i}(x,w)$ being a random variable with expectation zero. We handle each term
separately.

%
\paragraph{Handling $Q_1$.} Define the random variable
\[
\mathbf{Q}_1 = \E_{x \sim X,w \sim [D]}[Q_{1}(x,w)]\,.
\]
Unfortunately, the random variables $Q_{1}(x,w)$, 
for $x \in X$ and
$w \in [D]$, are not independent. In particular, it may be the case that $Q_1(x,w_1)$
and $Q_1(x,w_2)$ query the same input to $\ext$.
The next observation will help us overcome
this issue.
\begin{claim}\label{claim:partition}
Assume there exists a subset $V \subseteq [D]$ such that 
$\set{\adv_{i}(w)}_{i \in [t],w \in V}$ are all distinct, and enumerate
$\set{Q_1(x,w)}_{x \in X,w \in V} = Q_{1},\ldots,Q_{s = K|V|}$ arbitrarily. Then, for every $i \in [s]$,
\[\E[Q_{i} \mid Q_{1},\ldots,Q_{i-1}]=0\,.\]
\end{claim}
\begin{proof}
Fix $i \in [s]$.
First note that for $x_a \neq x_b$, $Q_{a} = Q_1(x_a,w_a)$ and $Q_{b} = Q_1(x_b,w_b)$ 
are independent. Assume that $Q_i = Q_{1}(x,w)$
and let $\set{i_1,\ldots,i_\ell} \subseteq [i-1]$ be the indices that correspond to the same $x$,
\ie, each $Q_{i_j} = Q_{1}(x,w_j)$ for some $w_j \neq w$. Thus,
\[\E[Q_{i} \mid Q_{1},\ldots,Q_{i-1}]=\E[Q_{i} \mid Q_{i_1},\ldots,Q_{i_\ell}]\,.\]
Next, fix all values of $\ext$ queried by the $Q_{i_j}$-s. Keeping the notation
$Q_{i_j} = Q_1(x,w_j)$, this means we fix every
$\ext(x,\adv_{r}(w_j))$ for $j \in [\ell]$ and $r \in [t]$. These fixings do not
affect $Q_i$, by our assumption on $V$. Thus, under these fixings, $\E[Q_i] = 0$, 
as desired.
\end{proof}

We now argue that we can partition $[D]$ to a bounded number of such $V$-s.
\begin{lemma}
There exists a partition $[D] = V_1 \cup \cdots \cup V_L$ for $L = O(dt^2  2^{\Delta})$
such that for every $i \in [L]$, $\set{A_{j}(w)}_{j \in [t],w \in V_i}$ are all distinct.
\end{lemma}
\begin{proof}
Let $G' = (W = [D],U = [D],E_0)$ be the bipartite graph in which each $w \in W$ is
connected to $\adv_{1}(w),\ldots,\adv_{t}(w)$. Thus, $G$ is left-regular with degree $t$
and its right-degree is bounded by $t \cdot 2^{\Delta}$. Let $G = (V = [D],E)$ be the
two-step walk graph of $G'$. Namely, $(x,y) \in E$ if and only if there exists a path
$x \sim z \sim y$ in $G'$, where $x,y \in W$ and $z \in U$. 
Note that the maximal degree in $G$ is at most $t^{2} \cdot 2^{\Delta}$.
We will repeatedly use
the following standard claim, which can be shown by a simple greedy algorithm.
\begin{claim}\label{claim:aux}
Let $G$ be an undirected graph over $n$ vertices with maximal degree $\delta$.
Then, the size of the largest independent set in $G$ is at least 
%
%
$n/(\delta+1)$.
%
%
%
\end{claim}
The crucial observation is that an independent set in $G$ corresponds to a valid partition.
To see this, take any $w_1,w_2 \in V$ such that $(w_1,w_2) \notin E$. By definition, there are no
$r_1,r_2 \in [t]$ such that $\adv_{r_1}(w_1) = \adv_{r_2}(w_2)$. In light of this observation,
we can greedily define $V_1,\ldots,V_{L}$ in the following manner.
\begin{enumerate}
\item	Set $G_0 \leftarrow G$, $i \leftarrow 0$ and $\delta = t^{2}2^{\Delta}$.
\item	As long as $G_i$ has more than $2\delta$ vertices,
\begin{itemize}
\item	Let $V_{i+1}$ be the largest independent set in $G_{i}$.
\item	Remove $V_{i+1}$ and all its adjacent edges and denote the resulting graph by $G_{i+1}$.
Set $i \leftarrow i+1$.
\end{itemize} 
\item	The graph $G_i$ has $b \le 2\delta$ vertices. Put each of these vertices in a separate
set.
\item	The resulting partition is $V_{1},\ldots,V_{i},\ldots,V_{L = i+b}$.
\end{enumerate}
\expref{Claim}{claim:aux} guarantees that at each iteration, $V_{i+1}$ contains at least 
%
%
$1-1/(\delta+1)$
%
fraction of the remaining vertices.
Let $j$ be the smallest integer for which
\[
\left(1 - \frac{1}{\delta+1} \right)^{j} \cdot 2^{d} \le 2\delta\,.
\]
One can verify that $j = O(\delta d)$, so overall $L \le j+2\delta = O\left(dt^{2}2^{\Delta} \right)$, as desired.
\end{proof}
In light of the above lemma and \expref{Claim}{claim:partition}, we can define, for each $i \in [L]$, 
\[\mathbf{S}_i = \sum_{x\in X,w \in V_i}Q_{1}(x,w)\,,\]
so $\mathbf{Q}_1 = \frac{1}{KD}\sum_{i \in [L]}\mathbf{S}_i$. Note that every sequence in
$\mathbf{S}_i$ is a martingale, and also, that $|Q_{1}(x,w)| \le 1$ with probability $1$. Thus,
using Azuma's inequality,
\begin{align}\label{eq:non-explicit1}
\Pr\left[\left|\mathbf{Q}_1\right| > \frac{\eps}{2}\right] &= \Pr\left[ \left|\sum_{i \in [L]}\mathbf{S}_i\right| > \frac{\eps}{2}KD \right] \le \sum_{i \in [L]}\Pr\left[\left|\mathbf{S}_i\right| > \frac{\eps}{2L}KD\right] \nonumber \\
&\le \sum_{i \in [L]} 2\exp\left( -\frac{\left(\frac{\eps}{2L}KD\right)^2}{2K|V_i|} \right) \le 
2L \cdot \eee^{-\frac{KD}{8L^2}\eps^2}\,.
\end{align}

\paragraph*{Handling $Q_2$.} Similarly, we define
$\mathbf{Q}_2 = \E_{x \sim X,w \sim [D]}[Q_{2}(x,w)]$, but we write it
as 
\[
\mathbf{Q}_2 = \E_{x \sim X,y \sim Y}\left[\E_{w \sim [D]}\left[ T_{w}(Z(x,y)) \right] - \mu(T)\right]\,,
\]
\ie, we switched the order of $w$ and $y$. Now, define
\[
Q_{2}'(x,y) = \E_{w \sim [D]}\left[ T_{w}(Z(x,y)) \right] - \mu(T)\,,
\]
so $\mathbf{Q}_2 = \E_{x \sim X,y \sim [D]}[Q'_{2}(x,y)]$. We follow the same reasoning as before:
For arbitrary $y_1 \neq y_2$, $Q_{2}'(x,y_1)$ may depend on $Q_{2}'(x,y_2)$, but with the same
partitioning we can overcome the dependencies. Also, for any $x \in \B^n$ and $y \in [D]$ 
it holds that
\[
\E\left[ \E_{w \sim [D]}\left[ T_{w}(Z(x,y)) \right] \right] = \mu(T)\,,
\]
so overall,
\begin{equation}\label{eq:non-explicit2}
\Pr\left[\left|\mathbf{Q}_2\right| > \frac{\eps}{2}\right] \le 2L \cdot \eee^{-\frac{KD}{8L^2}\eps^2}
\end{equation}
as well.

%
\paragraph*{Putting it all together.}   %
Combining \expeqref{Equation}{eq:non-explicit1} and \expeqref{Equation}{eq:non-explicit2}, we get
\[
\Pr\left[\E_{x \sim X,w \sim [D]}\left[ Q(x,w) \right] > \eps\right] \le 4L \cdot
\eee^{-\frac{KD}{8L^2}\eps^2}
\le 2^{-c_{1}\frac{KD}{d^{2}t^{4}2^{2\Delta}}\eps^{2}+\log d + 2\log t + c_2}
\]
for some universal constants $c_1,c_2 > 0$. To complete our analysis, we require $\ext$
to work for any $X$, $\adv_1,\ldots,\adv_t$ and $T$. By the union bound, the probability
for a random $\ext$ to fail, denote it by $p$, is at given by
\begin{align*}
p &\le \binom{N}{K}D^{tD}2^{DM^t}2^{-c_{1}\frac{KD}{d^{2}t^{4}2^{2\Delta}}\eps^{2}+\log d + 2\log t + c_2} \\
&\le 2^{K\log\left( \frac{Ne}{K} \right) + tDd + DM^t - c_{1}\frac{KD}{d^{2}t^{4}2^{2\Delta}}\eps^{2}+\log d + 2\log t + c_2} \\
&\le 2^{K(n-k+2)+tDd+DM^t +\log d +2\log t + c_2 - c_{1}\frac{KD}{d^{2}t^{4}2^{2\Delta}}\eps^{2}}\,.
\end{align*}
To prove that $p < 1$ (in fact, we will show that $p \ll 1$) it is sufficient to argue that:
\begin{enumerate}
\item\label{eq:prob1}  $K(n-k+2) \le \frac{c_1}{4}\frac{KD}{d^{2}t^{4}2^{2\Delta}}\eps^{2}$, and,
\item\label{eq:prob2}  $D(td+M^t) +\log d +2\log t + c_2  \le \frac{c_1}{4}\frac{KD}{d^{2}t^{4}2^{2\Delta}}\eps^{2}$, or, $D(4td+M^t) \le \frac{c_1}{4}\frac{KD}{d^{2}t^{4}2^{2\Delta}}\eps^{2}$.
\end{enumerate}
\expeqref{Item}{eq:prob1} is true whenever 
\[
D \ge \frac{4}{c_1} \cdot \frac{(n-k+2)d^{2}t^{4}2^{2\Delta}}{\eps^2}\,.
\]
\expeqref{Item}{eq:prob2} it true whenever
\[
K \ge \frac{4}{c_1} \cdot \frac{(4td+M^t)d^{2}t^{4}2^{2\Delta}}{\eps^2}\,.
\]
The bounds on $d$ and $k$ follow from the above two inequalities.
\end{proof}






%% Theory of Computing
%% File: v021a008.bbl
%% Publication date: October 28, 2025
%% Authors: Gil Cohen, Dean Doron, and Shahar Samocha
\newcommand{\prelim}{Preliminary version in\ }\newcommand{\prelims}{Preliminary
  versions in\ }
% Define empty bibhead if not already defined
\providecommand{\bibhead}[1]{}
% Define tocrefpdfbookmark if not already defined
\expandafter\ifx\csname pdfbookmark\endcsname\relax%
  \providecommand{\tocrefpdfbookmark}{}
\else\providecommand{\tocrefpdfbookmark}{%
   \hypertarget{tocreferences}{}%
   \pdfbookmark[1]{References}{tocreferences}}%
\fi

\tocrefpdfbookmark
\begin{thebibliography}{10}

\bibitem{BKSSW10}\bibhead{BKSSW10}
{\sc Boaz Barak, Guy Kindler, Ronen Shaltiel, Benny Sudakov, and Avi
  Wigderson}: Simulating independence: {N}ew constructions of condensers,
  {R}amsey graphs, dispersers, and extractors.
\newblock {\em J. ACM}, 57(4):20:1--52, 2010.
\newblock [\epfmtdoi{10.1145/1734213.1734214}]

\bibitem{BRSW12}\bibhead{BRSW12}
{\sc Boaz Barak, Anup Rao, Ronen Shaltiel, and Avi Wigderson}: 2-source
  dispersers for {$n^{o(1)}$} entropy, and {R}amsey graphs beating the
  {F}rankl--{W}ilson construction.
\newblock {\em Ann. Math.}, 176(3):1483--1543, 2012.
\newblock [\epfmtdoi{10.4007/annals.2012.176.3.3}]

\bibitem{BDT20}\bibhead{BDT20}
{\sc Avraham Ben-Aroya, Dean Doron, and Amnon Ta-Shma}: Near-optimal erasure
  list-decodable codes.
\newblock In {\em Proc. 35th Comput. Complexity Conf. (CCC'20)}, volume 169,
  pp. 1:1--27. Schloss Dagstuhl--Leibniz-Zentrum fuer Informatik, 2020.
\newblock [\epfmtdoi{10.4230/LIPIcs.CCC.2020.1}]

\bibitem{BADTS17}\bibhead{BADTS17}
{\sc Avraham Ben-Aroya, Dean Doron, and Amnon Ta-Shma}: An efficient reduction
  from two-source to nonmalleable extractors: Achieving near-logarithmic
  min-entropy.
\newblock {\em SIAM J. Comput.}, 51(2):31--49, 2022.
\newblock [\epfmtdoi{10.1137/17M1133245}]

\bibitem{Bourgain05}\bibhead{Bourgain05}
{\sc Jean Bourgain}: More on the sum-product phenomenon in prime fields and its
  applications.
\newblock {\em Int. J. Number Theory}, 1(1):1--32, 2005.
\newblock [\epfmtdoi{10.1142/S1793042105000108}]

\bibitem{CGL16}\bibhead{CGL16}
{\sc Eshan Chattopadhyay, Vipul Goyal, and Xin Li}: Non-malleable extractors
  and codes, with their many tampered extensions.
\newblock In {\em Proc. 48th STOC}, pp. 285--298. ACM Press, 2016.
\newblock [\epfmtdoi{10.1145/2897518.2897547}]

\bibitem{CL16}\bibhead{CL16}
{\sc Eshan Chattopadhyay and Xin Li}: Explicit non-malleable extractors,
  multi-source extractors and almost optimal privacy amplification protocols.
\newblock In {\em Proc. 57th FOCS}, pp. 158--167. IEEE Comp. Soc., 2016.
\newblock [\epfmtdoi{10.1109/FOCS.2016.25}]

\bibitem{CZ19}\bibhead{CZ19}
{\sc Eshan Chattopadhyay and David Zuckerman}: Explicit two-source extractors
  and resilient functions.
\newblock {\em Ann. Math.}, 189(3):653--705, 2019.
\newblock [\epfmtdoi{10.4007/annals.2019.189.3.1}]

\bibitem{CG88}\bibhead{CG88}
{\sc Benny Chor and Oded Goldreich}: Unbiased bits from sources of weak
  randomness and probabilistic communication complexity.
\newblock {\em SIAM J. Comput.}, 17(2):230--261, 1988.
\newblock [\epfmtdoi{10.1137/0217015}]

\bibitem{Cohen16b}\bibhead{Cohen16b}
{\sc Gil Cohen}: Local correlation breakers and applications to three-source
  extractors and mergers.
\newblock {\em SIAM J. Comput.}, 45(4):1297--1338, 2016.
\newblock [\epfmtdoi{10.1137/15M1029837}]

\bibitem{Cohen16c}\bibhead{Cohen16c}
{\sc Gil Cohen}: Making the most of advice: new correlation breakers and their
  applications.
\newblock In {\em Proc. 57th FOCS}, pp. 188--196. IEEE Comp. Soc., 2016.
\newblock [\epfmtdoi{10.1109/FOCS.2016.28}]

\bibitem{Cohen16d}\bibhead{Cohen16d}
{\sc Gil Cohen}: Non-malleable extractors---{N}ew tools and improved
  constructions.
\newblock In {\em Proc. 31st Comput. Complexity Conf. (CCC'16)}, pp. 8:1--29.
  Schloss Dagstuhl--Leibniz-Zentrum fuer Informatik, 2016.
\newblock [\epfmtdoi{10.4230/LIPIcs.CCC.2016.8}]

\bibitem{Cohen17}\bibhead{Cohen17}
{\sc Gil Cohen}: Towards optimal two-source extractors and {R}amsey graphs.
\newblock In {\em Proc. 49th STOC}, pp. 1157--1170. ACM Press, 2017.
\newblock [\epfmtdoi{10.1145/3055399.3055429}]

\bibitem{Cohen16}\bibhead{Cohen16}
{\sc Gil Cohen}: Two-source dispersers for polylogarithmic entropy and improved
  {R}amsey graphs.
\newblock {\em SIAM J. Comput.}, 50(3):30--67, 2021.
\newblock \prelim \href{https://doi.org/10.1145/2897518.2897530}{STOC'16}.
\newblock [\epfmtdoi{10.1137/16M1096219}]

\bibitem{CRS12}\bibhead{CRS12}
{\sc Gil Cohen, Ran Raz, and Gil Segev}: Nonmalleable extractors with short
  seeds and applications to privacy amplification.
\newblock {\em SIAM J. Comput.}, 43(2):450--476, 2014.
\newblock \prelim \href{https://doi.org/10.1109/CCC.2012.21}{CCC'12}.
\newblock [\epfmtdoi{10.1137/130908634}]

\bibitem{CS16}\bibhead{CS16}
{\sc Gil Cohen and Leonard~J. Schulman}: Extractors for near logarithmic
  min-entropy.
\newblock In {\em Proc. 57th FOCS}, pp. 178--187. IEEE Comp. Soc., 2016.
\newblock [\epfmtdoi{10.1109/FOCS.2016.27}]

\bibitem{DLWTZ11}\bibhead{DLWTZ11}
{\sc Yevgeniy Dodis, Xin Li, Trevor~D. Wooley, and David Zuckerman}: Privacy
  amplification and non-malleable extractors via character sums.
\newblock In {\em Proc. 52nd FOCS}, pp. 668--677. IEEE Comp. Soc., 2011.
\newblock [\epfmtdoi{10.1109/FOCS.2011.67}]

\bibitem{DW09}\bibhead{DW09}
{\sc Yevgeniy Dodis and Daniel Wichs}: Non-malleable extractors and symmetric
  key cryptography from weak secrets.
\newblock In {\em Proc. 41st STOC}, pp. 601--610. ACM Press, 2009.
\newblock [\epfmtdoi{10.1145/1536414.1536496}]

\bibitem{DKSS13}\bibhead{DKSS13}
{\sc Zeev Dvir, Swastik Kopparty, Shubhangi Saraf, and Madhu Sudan}: Extensions
  to the method of multiplicities, with applications to {K}akeya sets and
  mergers.
\newblock {\em SIAM J. Comput.}, 42(6):2305--2328, 2013.
\newblock [\epfmtdoi{10.1137/100783704}]

\bibitem{GW21}\bibhead{GW21}
{\sc Oded Goldreich and Avi Wigderson}: Robustly self-ordered graphs:
  Constructions and applications to property testing.
\newblock {\em TheoretiCS}, 1(8788), 2022.
\newblock \prelim \href{https://doi.org/10.4230/LIPIcs.CCC.2021.12}{CCC'21}.
\newblock [\epfmtdoi{10.46298/theoretics.22.1}]

\bibitem{GUV09}\bibhead{GUV09}
{\sc Venkatesan Guruswami, Christopher Umans, and Salil Vadhan}: Unbalanced
  expanders and randomness extractors from {P}arvaresh--{V}ardy codes.
\newblock {\em J. ACM}, 56(4):20:1--34, 2009.
\newblock [\epfmtdoi{10.1145/1538902.1538904}]

\bibitem{Li15}\bibhead{Li15}
{\sc Xin Li}: Three-source extractors for polylogarithmic min-entropy.
\newblock In {\em Proc. 56th FOCS}, pp. 863--882. IEEE Comp. Soc., 2015.
\newblock [\epfmtdoi{10.1109/FOCS.2015.58}]

\bibitem{Li17}\bibhead{Li17}
{\sc Xin Li}: Improved non-malleable extractors, non-malleable codes and
  independent source extractors.
\newblock In {\em Proc. 49th STOC}, pp. 1144--1156. ACM Press, 2017.
\newblock [\epfmtdoi{10.1145/3055399.3055486}]

\bibitem{Li19}\bibhead{Li19}
{\sc Xin Li}: Non-malleable extractors and non-malleable codes: partially
  optimal constructions.
\newblock In {\em Proc. 34th Comput. Complexity Conf. (CCC'19)}, pp. 28:1--49.
  Schloss Dagstuhl--Leibniz-Zentrum fuer Informatik, 2019.
\newblock [\epfmtdoi{10.4230/LIPIcs.CCC.2019.28}]

\bibitem{LRVW03}\bibhead{LRVW03}
{\sc Chi-Jen Lu, Omer Reingold, Salil Vadhan, and Avi Wigderson}: Extractors:
  Optimal up to constant factors.
\newblock In {\em Proc. 35th STOC}, pp. 602--611. ACM Press, 2003.
\newblock [\epfmtdoi{10.1145/780542.780630}]

\bibitem{NZ96}\bibhead{NZ96}
{\sc Noam Nisan and David Zuckerman}: Randomness is linear in space.
\newblock {\em J. Comput. System Sci.}, 52(1):43--52, 1996.
\newblock [\epfmtdoi{10.1006/jcss.1996.0004}]

\bibitem{RTS00}\bibhead{RTS00}
{\sc Jaikumar Radhakrishnan and Amnon Ta-Shma}: Bounds for dispersers,
  extractors, and depth-two superconcentrators.
\newblock {\em SIAM J. Discr. Math.}, 13(1):2--24, 2000.
\newblock [\epfmtdoi{10.1137/S0895480197329508}]

\bibitem{Raz05}\bibhead{Raz05}
{\sc Ran Raz}: Extractors with weak random seeds.
\newblock In {\em Proc. 37th STOC}, pp. 11--20. ACM Press, 2005.
\newblock [\epfmtdoi{10.1145/1060590.1060593}]

\bibitem{Shaltiel11}\bibhead{Shaltiel11}
{\sc Ronen Shaltiel}: An introduction to randomness extractors.
\newblock In {\em Proc. 38th Internat. Colloq. on Automata, Languages, and
  Programming (ICALP'11)}, pp. 21--41. Springer, 2011.
\newblock [\epfmtdoi{10.1007/978-3-642-22012-8\_2}]

\bibitem{TSU12}\bibhead{TSU12}
{\sc Amnon Ta-Shma and Christopher Umans}: Better condensers and new extractors
  from {P}arvaresh--{V}ardy codes.
\newblock In {\em Proc. 27th IEEE Conf. on Comput. Complexity (CCC'12)}, pp.
  309--315. IEEE Comp. Soc., 2012.
\newblock [\epfmtdoi{10.1109/CCC.2012.25}]

\bibitem{Trevisan01}\bibhead{Trevisan01}
{\sc Luca Trevisan}: Extractors and pseudorandom generators.
\newblock {\em J. ACM}, 48(4):860--879, 2001.
\newblock [\epfmtdoi{10.1145/502090.502099}]

\bibitem{Vadhan12}\bibhead{Vadhan12}
{\sc Salil~P. Vadhan}: Pseudorandomness.
\newblock {\em Found. Trends Theor. Comp. Sci.}, 7(1--3):1--336, 2012.
\newblock [\epfmtdoi{10.1561/0400000010}]

\end{thebibliography}
\end{document}
